\subsection{Κατηγορίες ρύπων}


\subsubsection{Αέρινα σωματίδια}



%%%%%%%%%%%%%%%%%%%%%%%%%%%%%%%%%%%%%%%%%%%%%%%%
Τα αέρινα σωματίδια (PM) αποτελούνται από στερεά σωματίδια αιωρούμενα στον αέρα, τα οποία μπορούν να εισέλθουν στο αναπνευστικό σύστημα, κυρίως όταν η διάμετρός τους είναι μικρότερη από 2,5 μm. Στην ίδια κατηγορία εντάσσονται και ίνες, όπως το υαλοβάμβακας και η αμίαντος (Carazo Fernández et al., 2013; Leung, 2015; Royal College of Physicians, 2016). Τα σωματίδια από το εξωτερικό περιβάλλον εισέρχονται στους εσωτερικούς χώρους μέσω του συστήματος αερισμού, ενώ οι εσωτερικές πηγές περιλαμβάνουν συσκευές καύσης, όπως φούρνους, θερμαντικά σώματα ή κουζίνες, καθώς και καπνό τσιγάρου και τζάκια. Επιπλέον, τα PM μπορεί να προκύπτουν από αντιδράσεις όζοντος με ορισμένα πτητικά οργανικά συστατικά (VOC), όπως τα τερπένια (Carazo Fernández et al., 2013; European Environment Agency, 2017; Royal College of Physicians, 2016; SCHER, 2007). Η χρήση βιομάζας ως καυσίμου μπορεί να οδηγήσει σε επίπεδα εσωτερικών σωματιδίων υψηλότερα από αυτά που παρατηρούνται σε μολυσμένες πόλεις.

Η μακροχρόνια έκθεση σε αέρινα σωματίδια μπορεί να προκαλέσει πλήθος αναπνευστικών και καρδιαγγειακών προβλημάτων, από ερεθισμό ματιών, μύτης, λαιμού και βρόγχων, καθώς και άσθμα, έως ίνωση, ανθρακώση και καρκίνο των πνευμόνων (Carazo Fernández et al., 2013; Jantunen et al., 2011; Leung, 2015; SCHER, 2007; US Environmental Protection Agency, 1995). Στους εσωτερικούς χώρους, όπως κουζίνες και υπνοδωμάτια, οι συγκεντρώσεις PM2.5 κυμαίνονται μεταξύ 25–1526 μg/m³ και 13–27 μg/m³, αντίστοιχα, ενώ σε γραφεία και σχολεία έχουν καταγραφεί τιμές 10–44 μg/m³ και 3–23 μg/m³. Τα εμπορικά κέντρα παρουσιάζουν επίσης υψηλές συγκεντρώσεις PM2.5, της τάξης των 74–164 μg/m³ (Armendáriz-Arnez et al., 2010; Barmparesos et al., 2018; Broderick et al., 2017; Canha et al., 2017; Datta et al., 2017; Mandin et al., 2017; Pokhrel et al., 2015; Sidhu et al., 2017; Singh et al., 2016; Vasile et al., 2016; Yip et al., 2017).

Particulate matter consists of solid particles suspended in air able to enter the respiratory tract, which mainly occurs when particle size is below 2.5 μm. Fibers, such as fiber glass and asbestos, can also be included in this group (Carazo Fernández et al., 2013; Leung, 2015; Royal College of Physicians, 2016). Outdoor PM enters indoor environments through ventilation, while indoor sources include combustion-based appliances like ovens, heaters or stoves, and also tobacco smoke and fireplaces. PM can also be originated by reactions between ozone and some VOCs (e.g. terpenes) (Carazo Fernández et al., 2013; European Environment Agency, 2017; Royal College of Physicians, 2016; SCHER, 2007). When biomass is used as a fuel, indoor levels of PM can exceed those of polluted cities. Long-term exposure to particulate matter can provoke conditions from several respiratory and cardiovascular problems such as eye, nose, throat and bronchial irritation and asthma to fibrosis, anthracosis and lung cancer (Carazo Fernández et al., 2013; Jantunen et al., 2011; Leung, 2015; SCHER, 2007; US Environmental Protection Agency, 1995). Kitchens and bedrooms may exhibit PM2.5 concentrations in the range of 25–1526 and 13–27 μg/m3, while concentrations of 10–44 and 3–23 μg/m3 have been recorded in offices and schools, respectively. Shopping complexes also present high concentrations of PM2.5 (74–164 μg/m3) (Armendáriz-Arnez et al., 2010; Barmparesos et al., 2018; Broderick et al., 2017; Canha et al., 2017; Datta et al., 2017; Mandin et al., 2017; Pokhrel et al., 2015; Sidhu et al., 2017; Singh et al., 2016; Vasile et al., 2016; Yip et al., 2017).



\cite{GONZALEZMARTIN2021128376}

Τα αέρινα σωματίδια (PM) αποτελούν μείγμα στερεών και υγρών σωματιδίων που αιωρούνται στον αέρα, περιλαμβάνοντας οξέα, οργανικές ενώσεις, αιθάλη, μέταλλα, σκόνη και έδαφος. Η ρύπανση από σωματίδια μπορεί να ταξινομηθεί ανάλογα με το μέγεθός τους (διάμετρος), περιλαμβάνοντας τα PM10 (2,5–10 μm), PM2.5 (<2,5 μm) και PM1.0 (<1,0 μm) [61]. Τα PM10 επηρεάζουν τις ρινικές και στοματικές κοιλότητες, τον φάρυγγα, τον λάρυγγα και τον ανώτερο τράχηλο. Τα PM2.5 είναι λεπτά εισπνεόμενα σωματίδια που εναποτίθενται στην επιφάνεια των επιθηλιακών κυττάρων των βρογχιολίων και των κυψελίδων, ενώ τα PM1.0 μπορούν να διεισδύσουν σε εσωτερικά όργανα, όπως η καρδιά και ο εγκέφαλος [62,63]. Τα PM2.5 και PM1.0 μπορούν να προκαλέσουν πνευμονικές λοιμώξεις και να δημιουργήσουν αγγειακές και ενδοθηλιακές δυσλειτουργίες, διαταραχές της μεταβλητότητας της καρδιακής συχνότητας, διαταραχές της πήξης και της καρδιακής αυτόνομης λειτουργίας [64]. Η έκθεση σε PM εκτιμάται ότι προκαλεί περίπου 3,3 εκατομμύρια θανάτους παγκοσμίως ετησίως [65]. Παιδιά, ηλικιωμένοι και άτομα με καρδιοαναπνευστικά προβλήματα αποτελούν τις πληθυσμιακές ομάδες υψηλού κινδύνου για τους ρύπους PM [50].

PM is a mixture of solid and liquid particles embodied in the air, including acids, organic chemicals, soot, metals, soil, and dust. Particle pollution can be categorized by its size (diameter), which includes PM10 (2.5 µm to 10 µm), PM2.5 (less than 2.5 µm) and PM1.0 (less than 1.0 µm) [61]. PM10 affects the nasal and oral cavities, the pharynx, the larynx, and the upper trachea. PM2.5 are fine inhalable particles that form sediments on the surface of epithelial cells in the bronchioles and alveoli. PM1.0 can lead inward to internal organs, including the heart and brain [62,63]. “PM2.5 and PM1.0 can lead to pulmonary infection and generate vascular and endothelial dysfunction, alterations in heart rate variability, coagulation, and cardiac autonomic function” [64]. PM is estimated to cause of 3.3 million deaths per year worldwide [65]. Children, the elderly, and people with heart and lung disease are the high-risk populations for PM pollutants [50].

\cite{su12219045}

Τα αέρινα σωματίδια (PM) αποτελούν σύνθετο μείγμα οργανικών και ανόργανων χημικών ενώσεων, περιλαμβάνοντας οργανικό άνθρακα (OC) και στοιχειακό άνθρακα (EC), των οποίων οι συγκεντρώσεις διαφέρουν ανάλογα με την εποχή και την γεωγραφική θέση [16,17]. Τα εσωτερικά σωματίδια περιλαμβάνουν αέρια σωματίδια διαφορετικών κλασμάτων μεγέθους (συνήθως αναφερόμενα ως PM10, PM2.5 και υπερλεπτά σωματίδια – UFP), τα οποία εισέρχονται από το εξωτερικό περιβάλλον, καθώς και σωματίδια που παράγονται εντός του χώρου [18,19]. Οι κύριες εσωτερικές πηγές PM περιλαμβάνουν το κάπνισμα, το μαγείρεμα (ιδιαίτερα με χρήση κηροζίνης και καυσίμων βιομάζας), ξυλόσομπες και φούρνους, τη χρήση θυμιάματος και κεριών για τα λεπτά σωματίδια (PM2.5), καθώς και δραστηριότητες καθαρισμού, παρουσία κατοικιδίων και ανθρώπινες κινήσεις για τα χονδρόκοκκα σωματίδια [20–26].

Εβδομήντα τρεις μελέτες ανέφεραν συγκεντρώσεις εσωτερικών PM2.5, με μέσες τιμές από 1,7 μg/m³ στην πόλη του Κεμπέκ, Καναδάς [27], έως 428,6 μg/m³ σε σπίτια με χρήση σωλήνα νερού (hookah) στο Ντουμπάι, ΗΑΕ [28] (Πίνακας S1). Το μαγείρεμα και το κάπνισμα αποτελούν τις κύριες πηγές PM2.5 σε οικίες στο Ηνωμένο Βασίλειο [29,30] και στις ΗΠΑ [31]. Επιπλέον του εσωτερικού καπνίσματος [32,33], τα επίπεδα PM2.5 επηρεάζονται σημαντικά από τη χρήση άνθρακα για μαγείρεμα [34] και τις εκπομπές οχημάτων [35] στην Κίνα. Η χρήση κηροζίνης και βιομάζας για μαγείρεμα αύξησε τις συγκεντρώσεις PM2.5 εσωτερικών χώρων σε επίπεδα επτά φορές υψηλότερα από τις οδηγίες του Παγκόσμιου Οργανισμού Υγείας (WHO) στη Ντάκα, Μπαγκλαντές [36].

Τριάντα επτά μελέτες ανέφεραν συγκεντρώσεις εσωτερικών PM10, με μέσες τιμές από <11,0 μg/m³ στο Κεμπέκ, Καναδάς [37], έως 1275 μg/m³ σε σπίτια με χώμα στο δάπεδο στο Λάος [38]. Πολύ υψηλά επίπεδα PM10 αναφέρθηκαν επίσης σε Δελχί και Άγρα, Ινδία [39,40]. Οι σχετικά υψηλές συγκεντρώσεις PM10 σε εσωτερικούς χώρους στη νότια Ευρώπη σε σύγκριση με τη βόρεια και κεντρική Ευρώπη πιθανώς οφείλονται σε μεγαλύτερη εισροή μεταλλικής σκόνης από εξωτερικές πηγές [41]. Αντίστοιχα, υψηλότερα επίπεδα PM10 αναφέρθηκαν σε πορτογαλικές [41] και βρετανικές οικίες [42] σε σύγκριση με τα εξωτερικά περιβάλλοντα, αν και οι εξωτερικές πηγές παραμένουν σημαντικός συντελεστής των εσωτερικών συγκεντρώσεων. Η περιεκτικότητα σε OC ήταν υψηλότερη στους εσωτερικούς χώρους, υποδεικνύοντας την επιρροή των εσωτερικών πηγών όπως το κάπνισμα, το μαγείρεμα, η καύση βιομάζας και η κίνηση ανθρώπων, ενώ οι πηγές EC ήταν κυρίως εξωτερικές [41].

Οι συγκεντρώσεις εσωτερικών UFP επηρεάζονται από τα επίπεδα UFP εξωτερικών χώρων, τα χαρακτηριστικά του κτιρίου και τη διαρροή αέρα, τους ρυθμούς ανανέωσης αέρα (AER), τις εσωτερικές και εξωτερικές μετεωρολογικές παραμέτρους, την ανθρώπινη συμπεριφορά, καθώς και από εσωτερικές πηγές και δραστηριότητες, όπως το μαγείρεμα, η χρήση κεριών, οι θερμαντικές συσκευές, ο καπνός από τσιγάρο και ο εξοπλισμός γραφείου [43–45]. Σε μια μελέτη στην Καλιφόρνια, ΗΠΑ, το μαγείρεμα προκάλεσε τη μεγαλύτερη εσωτερική έκθεση σε UFP [46].

Particulate matter (PM) is a complex mixture of organic and inorganic chemicals, including organic carbon (OC) and elemental carbon (EC), which vary by season and geographic location [16,17]. Indoor particles comprised ambient particles of different size fractions (typically reported as PM10, PM2.5, and UFP) that infiltrated from outdoors, and particles that were generated indoors [18,19]. Major indoor PM sources included smoking, cooking (particularly using kerosene and biomass fuels), wood stoves and furnaces, use of incense and candles for fine particles (PM2.5), and cleaning, presence of pets, and people’s movements for coarse particles [20,21,22,23,24,25,26].
Seventy-three studies reported indoor PM2.5, with mean concentrations ranging between 1.7 μg/m3 in Quebec City, Canada [27], and 428.6 μg/m3 in homes with hookah (i.e., water pipe) smoking in Dubai, UAE [28] (Table S1). Cooking and smoking were major indoor sources of PM2.5 in homes in the UK [29,30] and USA [31]. In addition to indoor smoking [32,33], PM2.5 levels were strongly influenced by the use of coal for cooking [34] and motor vehicle emissions [35] in China. The use of kerosene and biomass fuels for cooking increased indoor PM2.5 concentrations to a level 7-fold greater than the World Health Organization (WHO) guidelines in Dhaka, Bangladesh [36].
Thirty-seven studies reported indoor PM10, with mean concentrations ranging between < 11.0 in Quebec City, Canada [37], and 1275 μg/m3 in houses with soil floor in Lao PDR [38]. Very high PM10 levels were also reported in Delhi and Agra, India [39,40]. Relatively high indoor PM10 in southern Europe compared to northern and central Europe was possibly due to higher ingress of mineral dust from outdoor sources [41]. Higher indoor PM10 was reported in Portuguese [41] and UK homes [42] compared to outdoor levels, although outdoor sources were a significant contributor to indoor concentrations. OC content of PM10 was higher indoors than outdoors showing the influence of indoor sources such as smoking, cooking, biomass burning and movement of people, while EC sources were mainly outdoors [41].
Indoor UFP concentrations were affected by outdoor UFP concentrations, building characteristics and infiltration, air exchange rates (AERs), indoor and outdoor meteorological parameters, personal behaviours, as well as indoor sources and human activities such as cooking, candle burning, heating devices, environmental tobacco smoke (ETS), and office equipment [43,44,45]. In a study in California, USA, cooking caused the highest indoor UFP exposures [46].


\cite{ijerph17238972}

Τα αέρινα σωματίδια (PM) ορίζονται ως ανθρακικά σωματίδια που συνδέονται με προσροφημένες οργανικές ενώσεις και δραστικά μέταλλα. Τα κύρια συστατικά των PM περιλαμβάνουν θειικά, νιτρικά, ενδοτοξίνες, πολυκυκλικούς αρωματικούς υδρογονάνθρακες και βαρέα μέταλλα όπως σίδηρο, νικέλιο, χαλκό, ψευδάργυρο και βανάδιο [7]. Ανάλογα με το μέγεθος των σωματιδίων, τα PM ταξινομούνται γενικά σε: (i) χονδρόκοκκα σωματίδια, PM10, με διάμετρο <10 μm, (ii) λεπτά σωματίδια, PM2.5, με διάμετρο <2,5 μm, και (iii) υπερλεπτά σωματίδια, PM0.1, με διάμετρο <0,1 μm. Τα PM αποτελούν σημαντικό κίνδυνο για την υγεία, καθώς είναι εισπνεύσιμα, επηρεάζουν τους πνεύμονες και την καρδιά και προκαλούν σοβαρές επιπτώσεις [26,27]. Έχει παρατηρηθεί ότι τα επίπεδα PM στους εσωτερικούς χώρους συχνά υπερβαίνουν εκείνα των εξωτερικών χώρων [25].

Οι εσωτερικές πηγές PM περιλαμβάνουν (i) σωματίδια που εισέρχονται από το εξωτερικό περιβάλλον και (ii) σωματίδια που παράγονται από ανθρώπινες δραστηριότητες εντός του χώρου. Το μαγείρεμα, οι δραστηριότητες καύσης ορυκτών καυσίμων, το κάπνισμα, η λειτουργία μηχανών και οι οικιακές ασχολίες αποτελούν τους κύριους λόγους για τη διασπορά των PM εντός κτιρίων. Σε σύγκριση με τα PM10 και PM2.5, τα PM0.1 που παράγονται από καύση ορυκτών καυσίμων αποτελούν μεγαλύτερη απειλή για την υγεία λόγω της ικανότητάς τους να διεισδύουν στις μικρές αεροφόρες οδούς και στις κυψελίδες [26,27].

Έρευνες σχετικά με τις συγκεντρώσεις των κυριότερων εσωτερικών ρύπων έχουν δείξει ότι το μαγείρεμα και το κάπνισμα αποτελούν τις μεγαλύτερες πηγές PM εσωτερικών χώρων, ενώ οι δραστηριότητες καθαρισμού συνεισφέρουν σε μικρότερο βαθμό [28]. Το κάπνισμα θεωρείται σημαντική πηγή PM2.5, με εκτιμώμενες αυξήσεις σε σπίτια καπνιστών μεταξύ 25 και 45 μg/m³, με μεγαλύτερες συγκεντρώσεις το χειμώνα σε σχέση με το καλοκαίρι [29]. Το μαγείρεμα παράγει εκατομμύρια σωματίδια (~10^6 σωματίδια/cm³) μέσω της καύσης λαδιού, ξύλου και τροφίμων, εκ των οποίων τα περισσότερα είναι υπερλεπτά σωματίδια [30,31]. Αυτά τα λεπτά σωματίδια μπορούν να διασπείρονται όχι μόνο στην κουζίνα αλλά και στο σαλόνι και άλλους χώρους του κτιρίου, προκαλώντας αρνητικές επιπτώσεις στην υγεία των ενοίκων [31,32]. Επιπλέον, άλλες καθημερινές ανθρώπινες δραστηριότητες, όπως το περπάτημα και η καθιστική χρήση επίπλων, μπορούν να ανασηκώσουν σκόνη και να συνεισφέρουν περίπου στο 25% της εσωτερικής συγκέντρωσης PM [29].

Συνολικά, έχει βρεθεί ότι οι πηγές ανθρώπινων δραστηριοτήτων έχουν ρυθμούς εκπομπής που κυμαίνονται από 0,03 έως 0,5 mg·min⁻¹ για PM2.5 και από 0,1 έως 1,4 mg·min⁻¹ για PM10 [33].

 3.1. Particulate Matters
PM is defined as carbonaceous particles in association with adsorbed organic chemicals and reactive metals. PM’s main components are sulfates, nitrates, endotoxin, polycyclic aromatic hydrocarbons, and heavy metals (iron, nickel, copper, zinc, and vanadium) [7]. Depending on the particle size, PM generally is classified into (i) coarse particles, PM10 of diameter <10 µm; (ii) fine particles, PM2.5 of diameter <2.5 µm; and (iii) ultrafine particles, PM0.1 of diameter <0.1 µm. PM is especially concerning, as it is sometimes inhalable, affecting the lungs and heart and causing serious health effects. It has been shown that indoor PM levels often exceed outdoor ones [25]. Indoor PM sources include (i) particles that migrate from the outdoor environment and (ii) particles generated by indoor activities. Cooking, fossil fuel combustion activities, smoking, machine operation, and residential hobbies are the main reasons why PM is distributed inside of buildings. Compared with PM10 and PM2.5, PM0.1 created by fossil fuel combustion represents a greater threat to health due to its penetrability into the small airways as well as alveoli [26,27]. According to research about the concentration of major indoor pollutants, it has been indicated that cooking and cigarette smoking are the largest sources of indoor air PM, whereas cleaning activities often have a lesser contribution to indoor PM [28]. Smoking is known as a major source of indoor PM2.5, with estimated increases in homes with smokers ranging from 25 to 45 µg /m3 and the concentration in winter is greater than in summer [29]. For cooking activity, it was shown that cooking activities enable the emission of millions of particles (~106 particles/cm3) through the burning of oil, wood, and food and most of them are ultra-fine particles [30,31]. In addition, these fine particles can distribute not only to the kitchen but also spread to the living room and other areas in the building, thereby causing adverse effects to the occupants’ health [31,32]. Meanwhile, other normal human activities, such as walking around and sitting on furniture, are likely to resuspend house dust and contribute to 25% of the indoor PM concentrations [29]. In summary, it has been found that the source strengths for human activities ranged from 0.03 to 0.5 mg.min−1 for PM2.5 and from 0.1 and 1.4 mg.min−1 for PM10 [33].

\cite{ijerph17082927}

\subsubsection{Πτητικες χημικές ουσίες (Volatile Organic Compounds)}

Οι πτητικές οργανικές ενώσεις (VOCs) αναγνωρίζονται ως αέρια που περιέχουν ποικίλα χημικά συστατικά και εκπέμπονται από υγρά ή στερεά [34]. Η φορμαλδεΰδη, ένα άχρωμο αέριο με καυστική οσμή, που απελευθερώνεται από πολλά υλικά κατασκευής, όπως μοριοσανίδες, κόντρα πλακέ και κόλλες, αποτελεί μία από τις πιο διαδεδομένες VOCs. Οι συγκεντρώσεις των VOCs σε εσωτερικούς χώρους είναι τουλάχιστον 10 φορές υψηλότερες από ό,τι στους εξωτερικούς χώρους, ανεξαρτήτως της τοποθεσίας του κτιρίου [34,35].

Γενικά, οι εσωτερικές VOCs προέρχονται από τέσσερις κύριες πηγές: (i) ανθρώπινες δραστηριότητες, όπως μαγείρεμα, κάπνισμα και χρήση προϊόντων καθαρισμού και προσωπικής φροντίδας, (ii) παραγωγή από εσωτερικές χημικές αντιδράσεις, (iii) διείσδυση εξωτερικού αέρα μέσω συστημάτων αερισμού και διαρροών, και (iv) απελευθέρωση από τα υλικά κατασκευής του κτιρίου [35–39]. Οι συγκεντρώσεις των VOCs επηρεάζονται από παράγοντες όπως ο ρυθμός ανανέωσης του αέρα, η ηλικία και το μέγεθος του κτιρίου, οι ανακαινίσεις, τα επίπεδα VOCs εξωτερικού αέρα και το άνοιγμα θυρών και παραθύρων [40]. Επιπλέον, έχει καταδειχθεί ότι κατά τη διάρκεια 90 λεπτών μαγειρέματος εντοπίζονται περίπου 50 διαφορετικές VOCs [35].

Δεδομένου ότι οι VOCs είναι οργανικές ενώσεις με χαμηλό σημείο βρασμού (Tb) και εξατμίζονται εύκολα ακόμη και σε θερμοκρασία δωματίου, ο Παγκόσμιος Οργανισμός Υγείας (WHO) τις ταξινόμησε σε τέσσερις ομάδες: (i) πολύ πτητικές οργανικές ενώσεις (VVOCs) με Tb 50–100 °C, (ii) πτητικές οργανικές ενώσεις (VOCs) με 100 °C < Tb < 240 °C, (iii) ημι-πτητικές ανόργανες ενώσεις (SVOCs) με 240 °C < Tb < 380 °C και (iv) σωματιδιακή οργανική ύλη (POM) με Tb > 380 °C [41,42].

Η έκθεση στις VOCs που απελευθερώνονται από καταναλωτικά προϊόντα γίνεται κυρίως μέσω τριών οδών: εισπνοή, κατάποση ή δερματική επαφή. Οι περισσότεροι άνθρωποι δεν επηρεάζονται σοβαρά από βραχυχρόνια έκθεση σε χαμηλές συγκεντρώσεις VOCs, αλλά σε περιπτώσεις μακροχρόνιας έκθεσης, ορισμένες VOCs θεωρούνται επικίνδυνες για την ανθρώπινη υγεία, με πιθανή καρκινογόνο δράση [43]. Όσον αφορά τις SVOCs, η δερματική απορρόφηση απευθείας από τον αέρα συμβάλλει περισσότερο στην έκθεση σε σύγκριση με την εισπνοή [44,45].

Volatile organic compounds (VOCs) are recognized as gases containing a variety of chemicals emitted from liquids or solids [34]. Formaldehyde, a colorless gas with an acrid smell and which is released from many building materials, such as particleboard, plywood, and glues, is one of the most widespread VOCs. VOC concentrations in indoor environments are at least 10 times higher than outdoors, regardless of the building location [34,35]. Generally, indoor VOCs are generated from four main sources: (i) Human activities, including cooking, smoking, and the use of cleaning and personal care products; (ii) generation from indoor chemical reactions; (iii) penetration of outdoor air through infiltration and ventilation systems; and (iv) originating from building materials [35,36,37,38,39]. The VOCs concentration is able to be affected by air exchange rates, house age and size, building renovations, outdoor VOC levels, and door and window opening [40]. Moreover, it has been demonstrated that about 50 different VOCs are identified during 90-min cooking periods [35]. Because VOCs are organic chemicals that possess a low boiling point (Tb) and are easily volatilized even at room temperature, the WHO classified them into four groups: (i) Very volatile organic compounds (VVOCs) with Tb: 50–100 °C; (ii) volatile organic compounds (VOCs) with 100 °C < Tb < 240 °C; (iii) semi-volatile inorganic compounds (SVOCs) with 240 °C < Tb < 380 °C; and (iv) particulate organic matter (POM) with Tb > 380 °C [41,42]. Normally, exposure to VOCs released from consumer products is incurred via three main pathways: Inhalation, ingestion, or dermal contact. Most people are not seriously affected by short-term exposure to low concentrations of VOCs, but in cases of long-term exposure, some VOCs are considered to be harmful risks to human health, potentially causing cancer [43]. As for SVOCs, transdermal uptake directly from air has a higher contribution compared with intake via inhalation [44,45].

\cite{ijerph17082927}

Οι πτητικές οργανικές ενώσεις (VOCs) αποτελούν μια ποικιλία επικίνδυνων οργανικών χημικών ουσιών που συμμετέχουν σε ατμοσφαιρικές φωτοχημικές αντιδράσεις και θεωρούνται από τους κύριους παράγοντες που συμβάλλουν στο Σύνδρομο Κτιριακής Ασθένειας (SBS) [6,66,67]. Ο Παγκόσμιος Οργανισμός Υγείας (WHO) ταξινομεί τις VOCs, τόσο εσωτερικού όσο και εξωτερικού χώρου, σε τέσσερις κατηγορίες με βάση τα σημεία βρασμού τους: πολύ πτητικές οργανικές ενώσεις (VVOCs), VOCs και ημι-πτητικές οργανικές ενώσεις (SVOCs) [68].

Πολλές μελέτες έχουν δείξει ότι οι συγκεντρώσεις πολλών VOCs σε εσωτερικούς χώρους είναι σημαντικά υψηλότερες από εκείνες των εξωτερικών χώρων [69–71]. Οι κύριες πηγές εσωτερικών VOCs περιλαμβάνουν υλικά κατασκευής με υψηλές εκπομπές, έπιπλα, αερολύματα, φυτοφάρμακα, υφάσματα, καθώς και εξοπλισμό γραφείου όπως φωτοτυπικά μηχανήματα και εκτυπωτές λέιζερ [6,67,70,72]. Η Υπηρεσία Προστασίας Περιβάλλοντος των ΗΠΑ (US EPA) έχει εκδώσει κατάλογο επικίνδυνων ατμοσφαιρικών ρύπων, που περιλαμβάνει συνολικά 187 VOCs [73]. Επιπλέον, το πρότυπο ANSI/ASHRAE 62.1-2016 παρέχει τις Επίπεδα Αναφοράς Έκθεσης (Reference Exposure Levels, RELs) για 32 συγκεκριμένες VOCs εσωτερικού χώρου για τον γενικό πληθυσμό [74].

Οι πιο συνηθισμένες εσωτερικές VOCs—όπως το βενζόλιο, η αιθυλένιο, η φορμαλδεΰδη, η μεθυλενιοχλωρίδιο, η τετραχλωροαιθυλένιο, η τολουόλη, η ξυλόλη και η 1,3-βουταδιένιο—έχουν αποδειχθεί ότι αποτελούν καρκινογόνες, ερεθιστικές και τοξικές ουσίες για τον άνθρωπο [75–77]. Τα TVOCs χρησιμοποιούνται ως μέτρο της συνολικής συγκέντρωσης VOCs σε εσωτερικούς χώρους [78,79]. Η οξεία έκθεση σε TVOCs μπορεί να προκαλέσει ερεθισμό ματιών, μύτης και λαιμού, πονοκεφάλους, απώλεια συντονισμού και ναυτία, βλάβες στο ήπαρ, τα νεφρά και το κεντρικό νευρικό σύστημα, αναπνευστικά προβλήματα, ενώ ορισμένα μπορεί να προκαλέσουν καρκίνο [67]. Τα άτομα με άσθμα, τα μικρά παιδιά και οι ηλικιωμένοι είναι ιδιαίτερα ευάλωτοι στις επιπτώσεις της έκθεσης σε TVOCs [6,77,79,80].

VOCs represent a diverse set of hazardous organic chemicals that participate in atmospheric photochemical reactions, which are considered to be one of the major contributors to SBS [6,66,67]. The WHO classifies both indoor and outdoor VOCs as Very-VOCs (VVOCs), VOCs, and Semi-VOCs (SVOCs) according to their boiling points [68]. Many studies have shown that the concentrations of many indoor VOCs were markedly higher than their outdoor counterparts [69,70,71]. The main indoor VOC sources include high-emission building materials, furnishings, aerosol sprays, pesticides, dry knitted products, office equipment such as copiers, and laser printers [6,67,70,72]. The US EPA issued a list of hazardous air pollutants, which include a total of 187 VOCs [73]. In addition, the ANSI/ASHRAE 62.1-2016 standard provides the Reference Exposure Levels (RELs) of 32 specific types of indoor VOCs for the general population [74]. The most common indoor VOCs—such as benzene, ethylene, formaldehyde, methylene chloride, tetrachloroethylene, toluene, xylene, and 1,3-butadiene—have been proven to be contributors of human carcinogens, irritants and toxicants [75,76,77]. TVOCs are used as a measure of the total volume of indoor VOC concentrations [78,79]. “Acute exposure to indoor TVOCs can cause eye, nose and throat irritation, headaches, loss of coordination and nausea, damage to the liver, kidney and central nervous system, respiratory disease and some cause cancer” [67]. Asthmatics, young children, and elderly people are more vulnerable to the effects of exposure to TVOCs [6,77,79,80]. 

\cite{su12219045}

\subsubsection{Διαφορες ενώσεις}

\paragraph{3.3.1. Επίπεδα Έκθεσης, Πηγές και Καθοριστικοί Παράγοντες}
Οι πτητικές οργανικές ενώσεις (VOCs) απελευθερώνονται από ένα ευρύ φάσμα εσωτερικών και εξωτερικών πηγών μέσω διαδικασιών καύσης και εξάτμισης, όπως το κάπνισμα, οι εκπομπές από διαλύτες, οι ανακαινίσεις, τα οικιακά προϊόντα και τα φυτοφάρμακα [5,68]. Τυπικές VOCs που ανιχνεύονται σε εσωτερικούς χώρους περιλαμβάνουν βενζόλιο, τολουόλη, αιθυλοβενζόλιο και ξυλόλες (BTEX), που προέρχονται από καύση καυσίμων και ανακαινίσεις κατοικιών· βενζόλιο και στυρένιο από το κάπνισμα· αλκάνια από φυσικό αέριο· 1,4-διχλωροβενζόλιο από εντομοαπωθητικά· α-πινένιο από ξυλο-based υλικά κατασκευής· και λιμονένιο από αρωματικά καθαριστικά και προϊόντα πλυντηρίου [69–72].

Σαράντα δύο μελέτες ανέφεραν εσωτερικές συγκεντρώσεις μιας ή περισσότερων VOCs (εξαιρουμένων μελετών που αφορούσαν μόνο φορμαλδεΰδη ή άλλες καρβονύλες). Οι συγκεντρώσεις VOCs σε κατοικίες εξαρτώνται από πολλούς παράγοντες, όπως η ένταση των πηγών εκπομπής, οι ρυθμοί αερισμού και το οξειδωτικό περιβάλλον του εσωτερικού χώρου, που αντανακλούν διαφορές στη χρήση χημικών προϊόντων, τον σχεδιασμό και τα υλικά του κτιρίου, τη συμπεριφορά των κατοίκων και την εποχή [32,73,74]. Οι συγκεντρώσεις VOCs ήταν γενικά υψηλότερες στους εσωτερικούς χώρους σε σχέση με τους εξωτερικούς, ιδίως για το βενζόλιο, ιδιαίτερα σε ψυχρότερες εποχές λόγω μειωμένου αερισμού και χρήσης θερμαντικών συσκευών πετρελαίου και αερίου [75–81]. Οι εσωτερικές πηγές κυριαρχούσαν για τις περισσότερες VOCs και ιδιαίτερα για λιμονένιο, α-πινένιο, εξανάλη, πεντάλη, ο-ξυλένιο και ν-δοδεκάνιο [76,82]. Η χρήση τεχνητών αρωματικών προϊόντων αέρα σχετιζόταν σημαντικά με τις συνολικές VOCs (TVOC), καθώς και με βενζόλιο, τολουόλη και αιθυλοβενζόλιο [78].

Τριάντα εννέα μελέτες ανέφεραν εσωτερικά επίπεδα βενζολίου και/ή τολουόλης, με συγκεντρώσεις που κυμαίνονταν από 0,6–3,0 µg/m³ στο Kaunas, Λιθουανία [75], έως 24,8 µg/m³ στο Perth, Αυστραλία [74], και 325,5 µg/m³ στο Sapporo, Ιαπωνία [83], αντίστοιχα. Προηγούμενες ανασκοπήσεις από τον Sarigiannis et al. (2011) ανέφεραν παρόμοιες, αν και πιο περιορισμένες, συγκεντρώσεις βενζολίου (1,2–17,0 µg/m³) και τολουόλης (4,3–86,2 µg/m³) σε κατοικίες εντός της Ευρωπαϊκής Ένωσης [5].

Σπίτια όπου υπήρχε παθητικό κάπνισμα (ETS) παρουσίαζαν υψηλότερες συγκεντρώσεις σχεδόν όλων των VOCs, συμπεριλαμβανομένου του βενζολίου [70,72,76,84,85]. Οι συγκεντρώσεις τολουόλης, αιθυλοβενζολίου και ξυλωλών σε εσωτερικούς χώρους επηρεάζονταν κυρίως από εξωτερικές πηγές όπως η κυκλοφορία δρόμων [75,77], αν και οι συγκεντρώσεις τολουόλης και ο-ξυλενολίου ήταν επίσης αυξημένες σε κατοικίες με κάπνισμα [85], ενώ τα επίπεδα τολουόλης επηρεάζονταν από την παρουσία χαλιών [76]. Αυξημένα επίπεδα βενζολίου και τολουόλης βρέθηκαν επίσης σε υπόγεια και γκαράζ κοντά σε κατοικίες, πιθανώς λόγω χρήσης και αποθήκευσης διαλυτών, βενζίνης και μηχανημάτων που λειτουργούν με βενζίνη [70,82].

Η χρήση υλικών εσωτερικών χώρων χαμηλών εκπομπών και μη απορροφητικών, καθώς και κλιματικές συνθήκες που ευνοούν τον φυσικό αερισμό, οδήγησαν σε μειωμένα επίπεδα βενζολίου, τολουόλης και ξυλώλης σε Αθήνα, Ελλάδα, και Σεούλ, Κορέα [77,78]. Ωστόσο, δύο μελέτες από τη βόρεια Ευρώπη ανέφεραν συγκεντρώσεις τολουόλης και ξυλώλης σε εσωτερικούς χώρους σημαντικά υψηλότερες από τα επίπεδα εξωτερικού χώρου [86,87]. Διεθνής μελέτη έδειξε ότι τα επίπεδα VOCs σε νέες κατοικίες μειώθηκαν δραματικά και πλησίασαν τις μέσες τιμές των παλαιότερων κατοικιών μετά από ένα έτος από την κατασκευή [88]. Αντίθετα, μελέτη από το Χονγκ Κονγκ [89] δεν βρήκε συσχέτιση μεταξύ συγκεντρώσεων VOCs (εκτός της φορμαλδεΰδης) και ηλικίας του κτιρίου.

Η πυκνότητα κατοίκων συσχετίστηκε θετικά με τις συγκεντρώσεις VOCs και BTEX σε Αυστραλία [71] και Καναδά [70], αντίστοιχα. Οι συγκεντρώσεις VOCs ήταν αρνητικά συσχετισμένες με τον αερισμό. Τα αλκάνια και τα αρωματικά συσχετίστηκαν με την εγγύτητα σε κύριους δρόμους. Οι συγκεντρώσεις VOCs σε αυστραλιανές κατοικίες ήταν χαμηλότερες σε σχέση με μελέτες στη Βόρεια Αμερική και Ευρώπη, πιθανώς λόγω της πιο «διαπερατής» φύσης των κατοικιών στην Αυστραλία [71].

\paragraph{3.3.2. Υγειονομικές Επιπτώσεις}
Πολλές VOCs ταξινομούνται ως γνωστά ή πιθανά καρκινογόνα (ιδιαίτερα το βενζόλιο, γνωστό καρκινογόνο για τον άνθρωπο κυρίως συσχετισμένο με λευχαιμία), ερεθιστικά και τοξικά, ενώ η μέτρηση των TVOCs μπορεί να υποεκτιμά τους κινδύνους που σχετίζονται με μεμονωμένες ενώσεις [5,70]. Οι περισσότερες VOCs έχουν επίσης αναφερθεί ως σημαντικοί παράγοντες κινδύνου για άσθμα [73,74,90–92], με τη μεγαλύτερη συσχέτιση για το βενζόλιο, ακολουθούμενο από αιθυλοβενζόλιο και τολουόλη. Ωστόσο, μετά από προσαρμογή για συγχυτικούς παράγοντες, δεν βρέθηκαν σημαντικές συσχετίσεις μεταξύ έκθεσης σε βενζόλιο σε κατοικίες και αναπνευστικής υγείας σε βρέφη στην Ισπανία [79], ούτε μεταξύ χαμηλού επιπέδου έκθεσης σε VOCs (εκτός του d-λιμονενίου) και άσθματος σε παιδιά στην Πορτογαλία [93]. Η έκθεση σε υψηλές συγκεντρώσεις VOCs κατά την βρεφική ηλικία αύξησε τον κίνδυνο ατοπικής δερματίτιδας σε Κορεάτικα παιδιά [94]. Η κατοικιακή έκθεση σε α-πινένιο συσχετίστηκε με συμπτώματα λαιμού και αναπνευστικά στην Ιαπωνία [83], αν και δεν βρέθηκαν σημαντικές επιπτώσεις του α-πινενίου στα συμπτώματα του SBS [95]. Συνολικά, η έμφαση στο SBS στις μελέτες αυτές ήταν περιορισμένη σε σχέση με την προηγούμενη ανασκόπηση για την ποιότητα εσωτερικού αέρα (IAQ) από τον Jones (1999) [9].

 3.3. Volatile Organic Compounds
3.3.1. Exposure Levels, Sources and Determinants
VOCs are emitted from a very wide range of indoor and outdoor sources through combustion and evaporation, e.g., cigarette smoking, solvent-related emissions, renovations, household products and pesticides [5,68]. Typical VOCs found in the indoor environment include benzene, toluene, ethylbenzene and xylenes (BTEX) from fuel combustion and evaporation, and house renovations; benzene and styrene from cigarette smoking; alkanes from natural gas; 1,4-dichlorobenzene from moth repellents; a-pinene from wood-based building materials; and limonene from fragranced household cleaning and laundry products (e.g., [69,70,71,72]).
Forty-two studies reported indoor concentrations of one or more VOCs (excluding studies that only reported formaldehyde or other carbonyls). VOC levels in homes depended on many factors, such as the strength of emission sources, ventilation rates, and the indoor oxidative environment, which reflected differences in chemical use, building design and materials, occupant behaviour, and season (e.g., [32,73,74]). Reported VOC concentrations were generally higher indoors than outdoors, including for benzene, particularly in colder seasons due to reduced ventilation and the use of oil and gas heaters [75,76,77,78,79,80,81]. Indoor sources were dominant for most VOCs and particularly for limonene, a-pinene, hexanal, pentanal, o-xylene, and n-dodecane [76,82]. Use of artificial air freshener was significantly associated with total VOC (TVOC), benzene, toluene and ethylbenzene [78].
Thirty-nine studies reported indoor levels of benzene and/or toluene, ranging from 0.6 and 3.0 µg/m3 in Kaunas, Lithuania [75], to 24.8 μg/m3 in Perth, Australia [74], and 325.5 µg/m3 in Sapporo, Japan [83], respectively (Table S1). Studies previously reviewed by Sarigiannis et al. (2011) reported similar, though narrower, ranges of benzene (1.2–17.0 μg/m3) and toluene (4.3–86.2 μg/m3) in homes within the European Union [5]
Homes where ETS was present had higher concentrations of almost all VOCs including benzene [70,72,76,84,85]. Indoor toluene, ethylbenzene and xylenes concentrations were mainly dominated by outdoor sources such as road traffic [75,77], although toluene and o-xylene concentrations were also elevated in smoking homes [85] and toluene levels were affected by the presence of carpets [76]. Elevated levels of benzene and toluene were also found in basements and garages adjacent to residences, probably due to the use and storage of solvents, petrol and petrol-powered equipment [70,82].
Use of low-emission and non-absorbent indoor materials, and climatic conditions that favour natural ventilation, resulted in reduced indoor benzene, toluene and xylene levels in Athens, Greece, and in Seoul, Korea [77,78]. However, two studies from northern Europe reported indoor concentrations of toluene and xylene significantly higher than outdoor levels [86,87]. An international study showed that levels of VOCs in new homes decreased dramatically and were close to the mean values for the older homes after one year from construction [88]. However, a study from Hong Kong [89] found no relationship between VOC concentrations (other than formaldehyde) and building age.
Occupant density was positively associated with indoor VOC and BTEX concentrations in Australia [71] and Canada [70], respectively. Indoor VOC concentrations were negatively correlated with ventilation. Indoor alkanes and aromatics were associated with proximity to major roads. Levels of VOCs in Australian dwellings were lower than those from studies in North America and Europe, probably due to the leakier nature of Australian dwellings [71].
3.3.2. Health Effects
Many VOCs are classified as known or possible carcinogens (benzene in particular is a known human carcinogen mainly associated with leukaemia risk), irritants and toxicants, and measurement of TVOCs may underestimate the risks associated with individual compounds [5,70]. Most VOCs have also been reported to be significant risk factors for asthma [73,74,90,91,92], with the strongest association with benzene followed by ethylbenzene and toluene. However, after adjustment for confounding factors, no significant associations were found between residential benzene exposure and respiratory health in infants in Spain [79], or between low-level exposure to VOCs (except for d-limonene) and asthma status in children in Portugal [93]. Exposure to high concentration of VOCs during infancy increased the risk of atopic dermatitis in Korean children [94]. Residential exposure to a-pinene was associated with throat and respiratory symptoms in Japan [83]. However, there were no significant effects of a-pinene on SBS symptoms [95]. Overall, there was less attention on the SBS in the studies reviewed here than in the earlier IAQ review by Jones (1999) [9].


\paragraph{Formaldehyde and Other Carbonyl Compounds}

Η φορμαλδεΰδη απελευθερώνεται κυρίως από δομικά υλικά (π.χ. μονωτικά υλικά και προϊόντα συμπιεσμένου ξύλου), οικιακά προϊόντα (π.χ. χρώματα, καθαριστικά, φυτοφάρμακα, κόλλες), παρκέ και χαλιά, κάπνισμα (αν και όχι κυρία πηγή) και μη αεριζόμενες συσκευές καύσης καυσίμου [5,55,89,96]. Οι καρβονύλες μπορούν επίσης να εμφανιστούν ως δευτερογενείς ρύποι, προκύπτοντας από τη χημική αντίδραση πρωτογενών ρύπων με όζον.

Τριάντα τρεις μελέτες ανέφεραν εσωτερικές συγκεντρώσεις φορμαλδεΰδης, κυμαινόμενες από 7,5 µg/m³ στο Quebec City, Καναδάς [37], έως 134 µg/m³ σε νέες κατοικίες κατά τον πρώτο χρόνο σε διάφορες πόλεις της Ιαπωνίας [88]. Προηγούμενες ανασκοπήσεις ανέφεραν πιο περιορισμένες συγκεντρώσεις (12,3–46,1 µg/m³) σε κατοικίες της Ευρωπαϊκής Ένωσης [5]. Οι διαφορές στις εσωτερικές συγκεντρώσεις οφείλονταν στην ηλικία των κτιρίων, τη γεωγραφία, τα δομικά υλικά, τα έπιπλα και τα οικιακά προϊόντα. Οι τιμές φορμαλδεΰδης ήταν γενικά υψηλότερες σε νεότερα κτίρια, ιδιαίτερα σε εκείνα με ξύλινα πλαίσια ή νέα/ανακαινισμένα έπιπλα [55,76,78,89,97].

Μελέτες στην Ιαπωνία εκτίμησαν ότι οι υψηλότερες συγκεντρώσεις εμφανίζονται 0–2 έτη μετά την ανακαίνιση, με γραμμική συσχέτιση συγκέντρωσης–χρόνου μετά την ανακαίνιση. Η φορμαλδεΰδη και το α-πινένιο από ξύλινα υλικά απαιτούν μεγαλύτερο χρόνο αερισμού σε νέες κατοικίες σε σχέση με άλλες VOCs [88]. Σε πόλεις της Ανατολικής Ασίας όπως το Χονγκ Κονγκ, τα επίπεδα φορμαλδεΰδης και στυρενίου ήταν υψηλότερα σε σχέση με Ιαπωνία, Κίνα, Κορέα και Ταϊβάν, αντανακλώντας τη μεγαλύτερη χρήση οικιακών προϊόντων που περιέχουν αυτές τις ενώσεις [89].

Οι κύριες εσωτερικές πηγές ήταν κυρίαρχες για φορμαλδεΰδη, ακεταλδεΰδη και ακετόνη [76,104]. Οι συγκεντρώσεις φορμαλδεΰδης ήταν υψηλότερες κατά τους καλοκαιρινούς μήνες, λόγω αυξημένων θερμοκρασιών και αυξημένων επιπέδων οξειδωτικών όπως το όζον και τα ακόρεστα οργανικά συστατικά (π.χ. τερπένια) [105]. Ωστόσο, σε ορισμένες περιοχές της Ιταλίας (Bari), οι συγκεντρώσεις φορμαλδεΰδης και ακεταλδεΰδης ήταν υψηλότερες τον χειμώνα σε σχέση με την άνοιξη-καλοκαίρι, πιθανώς λόγω ανοιχτών παραθύρων κατά τις καλές καιρικές συνθήκες [97].

Η φορμαλδεΰδη είναι καρκινογόνος για τον άνθρωπο, σχετιζόμενη με αυξημένο κίνδυνο καρκίνου ρινοφαρυγγικού και λευχαιμίας [1]. Αποτελεί επίσης ερεθιστικό του ανώτερου αναπνευστικού, προκαλώντας συμπτώματα όπως ερεθισμό ματιών, μύτης και λαιμού, συνήθως σε συνθήκες εσωτερικής έκθεσης. Η έκθεση σε χαμηλές συγκεντρώσεις για σύντομο χρονικό διάστημα μπορεί να επηρεάσει τη λειτουργία του δερματικού φραγμού σε ασθενείς με ατοπική δερματίτιδα. Η χαμηλού επιπέδου έκθεση σε φορμαλδεΰδη μπορεί να αυξήσει τον κίνδυνο αλλεργικής ευαισθητοποίησης σε κοινά αεροαλλεργιογόνα σε παιδιά [78,97]. Νεαρά παιδιά που εκτέθηκαν σε πρόσφατες εσωτερικές ανακαινίσεις έδειξαν αυξημένο κίνδυνο εκζέματος [106]. Η φορμαλδεΰδη σχετίστηκε επίσης με νευρολογικά συμπτώματα, όπως δυσκολία συγκέντρωσης, σε μελέτη στα ΗΑΕ [107]. Αυξημένα επίπεδα αλδεϋδών εντός κατοικιών αυξάνουν τον πιθανό κίνδυνο για σύνδρομο Sick Building Syndrome (SBS) σε κατοίκους νέων κατοικιών στην Ιαπωνία [108].

3.4. Formaldehyde and Other Carbonyls
3.4.1. Exposure Levels, Sources and Determinants
Formaldehyde is mainly emitted from building materials (e.g., insulating materials and pressed-wood products), household products (e.g., paints, cleaning products, pesticides, adhesives), parquet flooring and carpets, smoking (although not a dominant source), and unvented fuel-burning appliances (e.g., [5,55,89,96]). Carbonyl compounds can also occur in the indoor environment as secondary pollutants as products of the reaction of a primarily pollutants with ozone.
Thirty-three studies reported indoor formaldehyde levels with concentrations ranging from 7.5 µg/m3 in Quebec City, Canada [37], to 134 µg/m3 in new homes (first year) across various cities in Japan [88] (Table S1). Studies previously reviewed by Sarigiannis et al. (2011) reported a narrower range of formaldehyde concentrations (12.3–46.1 μg/m3) in homes within the European Union [5]. Differences among indoor concentrations of formaldehyde were due to differences in building ages, geography, building materials, furniture, and household products. Formaldehyde levels were generally higher in newer houses [55,78,89], particularly in those with wooden frames or furniture bought new or restored [76,97]. Maruo et al. (2010) estimated the relationship between formaldehyde levels and the age and temperature of homes in Japan [98]. They obtained the highest formaldehyde concentrations for apartments 0–2 years after their renovation, with a linear relationship between formaldehyde concentration and years after renovation. Formaldehyde and a-pinene related to wooden materials needed a longer flushing period than other VOCs in new homes [88].
Indoor formaldehyde and styrene levels in Hong Kong were higher than in other East Asian cities in Japan, China, Korea, Hong Kong and Taiwan, reflecting the higher prevalence of household products and materials containing these chemicals in Hong Kong [89]. An Italian study [97] reported lower indoor concentrations of formaldehyde than those reported in Japan [99] and France [96,100]. Formaldehyde levels ranged very widely in households in England [101] and Canada [102,103].
Indoor sources were dominant for formaldehyde, acetaldehyde and acetone [76,104]. Formaldehyde concentrations tended to be higher in summer when temperatures were higher. For example, Rancière et al. [105] found an increase in the concentrations of formaldehyde as a result of indoor chemistry involving oxidants such as ozone, which is present at higher concentrations in summer, and unsaturated organic compounds, such as terpenes, in homes in Paris. However, formaldehyde and acetaldehyde levels were significantly higher in winter than in spring-summer in dwellings in Bari, Italy [97], possibly because of windows kept open during good weather.
3.4.2. Health Effects
Formaldehyde is carcinogenic to humans, based on increased risk of nasopharyngeal cancer and leukaemia [1]. It is also an irritant of the upper respiratory tract with symptoms such as eye, nose and throat irritation commonly associated with indoor exposure. Exposure to low concentrations of formaldehyde for a short period of time influences the skin barrier function in patients with atopic dermatitis. Low level exposure to indoor formaldehyde may increase the risk of allergic sensitization to common aeroallergens in children [78], although this risk was low in households in south Italy [97]. Young children experiencing recent indoor renovation in German houses showed increased risk of eczema [106]. Formaldehyde was associated with neurological symptoms (difficulty concentrating) in a study in UAE [107]. Elevated levels of indoor aldehydes increased the possible risk of SBS in residents living in new houses in Japan [108].

\paragraph{Naphthalene}

Η ναφθαλίνη, που αποτελεί VOC και δακτυλιοειδή πολυκυκλική αρωματική υδρογονάνθρακα (PAH), είναι πανταχού παρούσα ρύπος και πολύ υψηλές συγκεντρώσεις εντοπίζονται σε εσωτερικούς χώρους όταν χρησιμοποιείται ως εντομοαπωθητικό ή αρωματικό [82,116]. Άλλες πηγές που επηρεάζουν τα επίπεδα εσωτερικών χώρων περιλαμβάνουν, σε μικρότερο βαθμό, το κάπνισμα και τις εκπομπές από οχήματα.

Δεκατέσσερις μελέτες ανέφεραν επίπεδα ναφθαλίνης σε κατοικίες, κυμαινόμενα από 0,12 µg/m³ στο Kaunas, Λιθουανία [75], έως 26,3 µg/m³ (και σε μία περίπτωση >1000 µg/m³) σε κατοικίες στο Michigan, ΗΠΑ [82]. Τυπικές συγκεντρώσεις σε σπίτια μη καπνιστών κυμαίνονται μεταξύ 0,18–1,7 µg/m³ [117]. Οι εξωτερικές συγκεντρώσεις σε αστικές περιοχές ήταν συνήθως χαμηλότερες από τις εσωτερικές [117]. Η ναφθαλίνη ήταν ο κυρίαρχος PAH σε αστικό υπόβαθρο και σε παρακείμενες οδούς, τόσο εσωτερικά όσο και εξωτερικά, στην Agra, Ινδία [118].


Μεγάλες ποσότητες ναφθαλίνης στον αέρα μπορούν να προκαλέσουν ερεθισμό ματιών και αναπνευστικού συστήματος. Ο ΠΟΥ (2010) όρισε ως μέση ετήσια οδηγία ποιότητας αέρα εσωτερικών χώρων το 10 µg/m³, βάσει βλαβών στο αναπνευστικό σύστημα, περιλαμβανομένων όγκων στο ανώτερο αναπνευστικό (αποδεικνύονται σε πειράματα σε ζώα) και αιμολυτικής αναιμίας στον άνθρωπο [1]. Συνολικά, η ναφθαλίνη παρουσιάζει κινδύνους για την υγεία σε υποσύνολα κατοικιών όπου γίνεται ακατάλληλη χρήση εντομοαπωθητικών και αρωματικών προϊόντων [116].
 3.6. Naphthalene
3.6.1. Exposure Levels, Sources and Determinants
Naphthalene (both a VOC and a 2-ring PAH) is a ubiquitous pollutant, and very high concentrations are sometimes encountered indoors when it is used as an insect repellent or deodorant [82,116]. Other sources that have an impact on indoor levels include (to lesser extents) cigarette smoking and motor vehicle emissions.
Fourteen studies reported naphthalene levels in domestic indoor environments, with concentrations ranging from 0.12 μg/m3 in homes in Kaunas, Lithuania [75], to 26.3 μg/m3 (and in one case exceeding 1000 μg/m3) in houses in Michigan, USA [82]. Jia and Batterman (2010) suggested typical naphthalene concentrations ranged from 0.18 to 1.7 μg/m3 in non-smokers’ homes [117], which were less variable than the indoor concentrations reported in this review (Table S1). Outdoor concentrations in urban areas reported by Jia and Batterman (2010) were typically lower than indoor levels [117]. Naphthalene was the dominant PAH at both urban background and roadside locations (indoors and outdoors) in Agra, India [118].
3.6.2. Health Effects
Large amounts of naphthalene in the air can irritate the eyes and respiratory system. The WHO (2010) established an annual average IAQ guideline (10 μg/m3) based on respiratory tract lesions, including tumors in the upper respiratory tract demonstrated in animal studies, and hemolytic anemia in humans [1]. Overall, naphthalene presents health risks in a subset of homes where inappropriate use of repellents and deodorants takes place [116].

\cite{ijerph17238972}

\paragraph{Benzene}
Το βενζόλιο είναι μια πτητική οργανική ένωση ευρείας χρήσης στη βιομηχανία. Οι εσωτερικές πηγές βενζολίου περιλαμβάνουν τον περιβαλλοντικό καπνό καπνού (ETS), συσκευές καύσης και δομικά υλικά, όπως πολυμερή έπιπλα, χαλιά, χρώματα, διαλύτες, διάφορα πλαστικά και ξύλινα έπιπλα (Campagnolo et al., 2017; Carazo Fernández et al., 2013; Jantunen et al., 1998; Kotzias et al., 2005; Royal College of Physicians, 2016; SCHER, 2007; Singh et al., 2016; World Health Organization and WHO Regional Office for Europe, 2010). Οι εκπομπές βενζολίου (και άλλων ρύπων) από τα υλικά αυτά μειώνονται με την πάροδο του χρόνου· ως εκ τούτου, τα επίπεδα είναι σημαντικά υψηλότερα σε νέα ή πρόσφατα ανακαινισμένα κτίρια.

Οι κίνδυνοι για την υγεία από οξεία δηλητηρίαση θεωρούνται χαμηλοί (π.χ. πονοκέφαλοι μετά από έκθεση 5 ωρών σε συγκεντρώσεις 150–500 mg/m³), ωστόσο το βενζόλιο είναι αποδεδειγμένα καρκινογόνο και γενοτοξικό σε συνθήκες χρόνιας έκθεσης. Επιπλέον, μπορεί να προκαλέσει αιματολογικές παθήσεις καθώς και νευρολογικές και αναπαραγωγικές διαταραχές (Carazo Fernández et al., 2013; Kotzias et al., 2005; Leung, 2015; World Health Organization and WHO Regional Office for Europe, 2010). Δεν συνιστάται ασφαλές επίπεδο έκθεσης στο βενζόλιο.

Οι υψηλότερες συγκεντρώσεις βενζολίου έχουν καταγραφεί σε εμπορικούς χώρους (2,5–48 μg/m³), ακολουθούμενες από γραφεία (1,4–5,5 μg/m³), κατοικίες (0,7–4,4 μg/m³) και σχολεία (0,5–3 μg/m³) (Broderick et al., 2017; Dodson et al., 2008; Du et al., 2015; Edwards et al., 2001; Geiss et al., 2011; Kotzias et al., 2009; Langer et al., 2015; Rösch et al., 2014; Singh et al., 2016; Xu et al., 2016; Zielinska et al., 2012).


Benzene is an organic volatile compound widely used in industry. Indoor sources of benzene include ETS, combustion devices and construction materials such as polymeric furnishings, carpets, paints, solvents, different plastics and wooden furniture (Campagnolo et al., 2017; Carazo Fernández et al., 2013; Jantunen et al., 1998; Kotzias et al., 2005; Royal College of Physicians, 2016; SCHER, 2007; Singh et al., 2016; World Health Organization and WHO Regional Office for Europe, 2010). Benzene (and other pollutants) emissions from these materials decay with time, thus levels are much higher in new or renovated buildings. Health risks for acute intoxication are low (e.g. headaches for 5-h exposure to 150–500 mg/m3), although benzene is a proven carcinogenic and genotoxic chemical under chronic exposure. Benzene can also induce blood illnesses and neurological and reproductive problems (Carazo Fernández et al., 2013; Kotzias et al., 2005; Leung, 2015; World Health Organization and WHO Regional Office for Europe, 2010). No safe level of exposure is recommended. The highest level of benzene has been recorded in shopping areas (2.5–48 μg/m3) followed by offices (1.4–5.5 μg/m3), homes (0.7–4.4 μg/m3) and schools (0.5–3 μg/m3) (Broderick et al., 2017; Dodson et al., 2008; Du et al., 2015; Edwards et al., 2001; Geiss et al., 2011; Kotzias et al., 2009; Langer et al., 2015; Rösch et al., 2014; Singh et al., 2016; Xu et al., 2016; Zielinska et al., 2012).


\paragraph{Toluene}

Η τολουόλη, η οποία απαντάται στη βενζίνη, τα χρώματα, τις ρητίνες, τα συγκολλητικά, τα καλλυντικά και τις επιστρώσεις, αποτελεί έναν από τους συχνότερους ρύπους του εσωτερικού αέρα και τη σημαντικότερη συνιστώσα της ομάδας BTEX (βενζόλιο, τολουόλη, αιθυλοβενζόλιο και ξυλόλες). Οι κύριες πηγές τολουόλης σε εσωτερικούς χώρους περιλαμβάνουν τη διείσδυση από τον εξωτερικό αέρα, τον καπνό καπνού, τις συσκευές καύσης και ένα ευρύ φάσμα οικιακών προϊόντων (Carazo Fernández et al., 2013; Jantunen et al., 1998; Kotzias et al., 2005; Royal College of Physicians, 2016; SCHER, 2007; Singh et al., 2016).

Η τολουόλη δεν έχει επιβεβαιωθεί ως καρκινογόνος ουσία και, ως εκ τούτου, οι μακροχρόνιες καρκινογενετικές επιπτώσεις της δεν έχουν τεκμηριωθεί επαρκώς. Παρ’ όλα αυτά, η έκθεση σε τολουόλη έχει συσχετιστεί με αυξημένο κίνδυνο εμφάνισης άσθματος και άλλων αναπνευστικών παθήσεων, ενώ τόσο η βραχυχρόνια όσο και η μακροχρόνια έκθεση μπορεί να επηρεάσει το κεντρικό νευρικό σύστημα (Kotzias et al., 2005).

Οι υψηλότερες συγκεντρώσεις τολουόλης έχουν καταγραφεί σε εμπορικά κέντρα (15–164 μg/m³), ακολουθούμενες από χώρους γραφείων (6–32 μg/m³) και κατοικίες (3–20 μg/m³). Η παρουσία τολουόλης σε σχολεία είναι περιορισμένη, με τυπικές συγκεντρώσεις της τάξης των 1,8 μg/m³. Οι συγκεντρώσεις είναι συνήθως υψηλότερες σε νέα ή πρόσφατα ανακαινισμένα κτίρια, γεγονός που αποδίδεται στην εκπομπή ρύπων από τα νεοεγκατεστημένα υλικά (Broderick et al., 2017; Dodson et al., 2008; Du et al., 2015; Edwards et al., 2001; Geiss et al., 2011; Langer et al., 2015; Mandin et al., 2017; Rösch et al., 2014; Singh et al., 2016; Xu et al., 2016; Zhong et al., 2017; Zielinska et al., 2012).

Toluene, present in gasoline, paints, resins, adhesives, cosmetics and coatings, is one of the most common indoor air pollutants and the most abundant among BTEX (benzene, toluene, ethylbenzene and xylene). Indoor sources include infiltration from outdoor air, tobacco smoke, combustion devices and a variety of household products (Carazo Fernández et al., 2013; Jantunen et al., 1998; Kotzias et al., 2005; Royal College of Physicians, 2016; SCHER, 2007; Singh et al., 2016). Toluene is not a confirmed carcinogenic substance, and therefore no long-term health effects have been studied yet. Exposure to toluene increases risks of developing asthma and other respiratory conditions, while short and long-term exposure can affect central nervous system (Kotzias et al., 2005). Shopping complexes present the highest toluene concentrations (15–164 μg/m3), followed by offices (6–32 μg/m3) and homes (3–20 μg/m3). The presence of toluene in schools is marginal (1.8 μg/m3). Concentrations are typically higher in new or recently renovated buildings as a result of the new installed materials (Broderick et al., 2017; Dodson et al., 2008; Du et al., 2015; Edwards et al., 2001; Geiss et al., 2011; Langer et al., 2015; Mandin et al., 2017; Rösch et al., 2014; Singh et al., 2016; Xu et al., 2016; Zhong et al., 2017; Zielinska et al., 2012).

\paragraph{Ethylbenzene}


Η αιθυλοβενζόλη είναι ένα αρωματικό, άχρωμο υγρό σε θερμοκρασία περιβάλλοντος. Οι πηγές ρύπανσης περιλαμβάνουν πλαστικά, χρώματα, συγκολλητικά και άλλα προϊόντα στα οποία η αιθυλοβενζόλη χρησιμοποιείται ως διαλύτης κατά τη διαδικασία παραγωγής. Επιπλέον, η αιθυλοβενζόλη εκπέμπεται και κατά τις διεργασίες καύσης (Campagnolo et al., 2017; Jantunen et al., 1998; Leung, 2015; Singh et al., 2016).

Οι οξείες επιδράσεις σε σχετικά υψηλές συγκεντρώσεις περιλαμβάνουν ερεθισμό των οφθαλμών και του φάρυγγα, καθώς και ζάλη. Μέχρι σήμερα δεν έχουν τεκμηριωθεί μακροχρόνιες επιπτώσεις στην ανθρώπινη υγεία. Ωστόσο, η αιθυλοβενζόλη έχει ταξινομηθεί ως δυνητικά καρκινογόνος για τον άνθρωπο από τον Διεθνή Οργανισμό Έρευνας για τον Καρκίνο (International Agency for Research on Cancer, IARC) (ATSDR, 2007).

Συγκεντρώσεις αιθυλοβενζόλης έως και 16 μg/m³ έχουν ανιχνευθεί σε χώρους γραφείων και εμπορικά κέντρα. Αντιθέτως, οι κατοικίες και τα σχολεία εμφανίζουν μέγιστες συγκεντρώσεις της τάξης των 8 και 0,2 μg/m³, αντίστοιχα (Broderick et al., 2017; Dodson et al., 2008; Edwards et al., 2001; Rösch et al., 2014; Singh et al., 2016; Xu et al., 2016; Zhong et al., 2017; Zielinska et al., 2012).

Ethylbenzene is an aromatic, colorless liquid at ambient temperature. Sources of pollution include plastics, paints, adhesives and other products where ethylbenzene is used as a solvent during manufacturing. Ethylbenzene is also emitted in combustion processes (Campagnolo et al., 2017; Jantunen et al., 1998; Leung, 2015; Singh et al., 2016). Acute effects at relatively high concentrations include eye and throat irritation and dizziness. No long-term effects have been described yet in humans. Ethylbenzene has been classified as a potential human carcinogen by the International Agency for Cancer Research (ATSDR, 2007). Concentrations of ethylbenzene up to 16 μg/m3 have been detected in offices and shopping complexes. On the other hand, homes and schools exhibit maximum concentrations of 8 and 0.2 μg/m3, respectively (Broderick et al., 2017; Dodson et al., 2008; Edwards et al., 2001; Rösch et al., 2014; Singh et al., 2016; Xu et al., 2016; Zhong et al., 2017; Zielinska et al., 2012).

\paragraph{Xylene}

Οι ξυλόλες απαντώνται σε τρία διαφορετικά ισομερή: ορθο-, μετα- και παρα- (o-, m- και p-, αντίστοιχα). Οι πηγές ξυλολών σε εσωτερικούς χώρους περιλαμβάνουν χρώματα, συγκολλητικές ουσίες, βαφές, πολυμερή, καθαριστικά προϊόντα και φαρμακευτικά σκευάσματα (Campagnolo et al., 2017; Jantunen et al., 1998; Kotzias et al., 2005; Leung, 2015; SCHER, 2007; Singh et al., 2016).

Η οξεία έκθεση σε ξυλόλες μπορεί να προκαλέσει ερεθισμό των οφθαλμών και του φάρυγγα, κεφαλαλγία και ναυτία. Η μακροχρόνια έκθεση έχει συσχετιστεί με διαταραχές του αναπνευστικού, του γαστρεντερικού και του κεντρικού νευρικού συστήματος, καθώς και με επιπτώσεις στους πνεύμονες, τους νεφρούς, την καρδιά και το αναπαραγωγικό σύστημα. Οι ξυλόλες θεωρούνται ύποπτες καρκινογόνες ουσίες και έχουν συσχετιστεί με αυξημένο κίνδυνο εμφάνισης λευχαιμίας, μη-Hodgkin λεμφώματος και καρκίνου του παχέος εντέρου/ορθού (Kotzias et al., 2005).

Οι εμπορικοί χώροι συγκαταλέγονται μεταξύ των πλέον επιβαρυμένων περιβαλλόντων, με συγκεντρώσεις που κυμαίνονται από 1,3 έως 74 μg/m³, ενώ σε χώρους γραφείων έχουν καταγραφεί επίπεδα 2,2–16 μg/m³. Αντιθέτως, οι μέγιστες συγκεντρώσεις ξυλολών που έχουν μετρηθεί σε κατοικίες και σχολεία είναι σημαντικά χαμηλότερες, της τάξης των 3,1 και 0,3 μg/m³ αντίστοιχα (Broderick et al., 2017; Dodson et al., 2008; Edwards et al., 2001; Geiss et al., 2011; Langer et al., 2015; Mandin et al., 2017; Singh et al., 2016; Vasile et al., 2016; Xu et al., 2016; Zhong et al., 2017; Zielinska et al., 2012).

Xylene is present as three different isomers: ortho, meta and para (o-, m- and p-, respectively). Indoor sources of xylenes include paints, adhesives, dyes, polymers, cleaning products and pharmaceuticals. (Campagnolo et al., 2017; Jantunen et al., 1998; Kotzias et al., 2005; Leung, 2015; SCHER, 2007; Singh et al., 2016). An acute exposure to xylene can cause eye and throat irritation, headache and nausea. Long-term exposure is associated with problems in the respiratory, gastrointestinal and central nervous systems, lungs, kidney, heart and reproductive system. Xylene is a suspected carcinogenic chemical, associated with an increased risk of leukemia, non-Hodgkin’s lymphoma and colon/rectum cancer (Kotzias et al., 2005). Shopping areas, with levels ranging from 1.3 to 74 μg/m3, and offices, with levels of 2.2–16 μg/m3, rank among the most polluted sites. Maximum xylene concentrations of 3.1 and 0.3 μg/m3 have been measured at homes and schools (Broderick et al., 2017; Dodson et al., 2008; Edwards et al., 2001; Geiss et al., 2011; Langer et al., 2015; Mandin et al., 2017; Singh et al., 2016; Vasile et al., 2016; Xu et al., 2016; Zhong et al., 2017; Zielinska et al., 2012).


\paragraph{Naphthalene}

Η ναφθαλίνη είναι ένα λευκό κρυσταλλικό στερεό, το οποίο χρησιμοποιείται ευρέως ως πρώτη ύλη στη χημική βιομηχανία. Πηγές ναφθαλίνης σε εσωτερικούς χώρους περιλαμβάνουν τον καπνό του τσιγάρου, ελαττωματικές συσκευές καύσης, ζιζανιοκτόνα, ελαστικά υλικά και εντομοαπωθητικά (συνήθως μέσω επιμόλυνσης των ρούχων), καθώς και τον εξωτερικό αέρα (Kotzias et al., 2005; Royal College of Physicians, 2016; World Health Organization and WHO Regional Office for Europe, 2010).

Η εισπνοή αποτελεί τον κύριο μηχανισμό πρόσληψης της ναφθαλίνης στον άνθρωπο, αν και έχουν αναφερθεί περιπτώσεις τυχαίας κατάποσης σφαιριδίων ναφθαλίνης. Η οξεία δηλητηρίαση από ναφθαλίνη μπορεί να προκαλέσει αιμολυτική αναιμία και καταρράκτη. Δεν έχουν αναφερθεί μακροχρόνιες δυσμενείς επιπτώσεις στην υγεία ούτε τεκμηριωμένα στοιχεία καρκινογένεσης ή γονοτοξικότητας στον άνθρωπο (Kotzias et al., 2005; World Health Organization and WHO Regional Office for Europe, 2010).

Συγκεντρώσεις ναφθαλίνης που κυμαίνονται από 3 έως 26 μg/m³ έχουν καταγραφεί σε κατοικίες (Du et al., 2015; Xu et al., 2016; Zhong et al., 2017).

Naphthalene is a white crystalline solid widely used as a raw material in the chemical industry. Indoor sources include tobacco smoke, defective combustion devices, herbicides, rubber materials and insect repellents (usually contaminating clothes), as well as outdoor air (Kotzias et al., 2005; Royal College of Physicians, 2016; World Health Organization and WHO Regional Office for Europe, 2010). Inhalation is the principal naphthalene intake mechanism in humans, although accidental ingestions of mothballs have been reported. Acute naphthalene intoxication can induce hemolytic anemia and cataracts. No long-term adverse health effects or evidence of human carcinogenicity or genotoxicity have been reported (Kotzias et al., 2005; World Health Organization and WHO Regional Office for Europe, 2010). Naphthalene concentrations ranging from 3 to 26 μg/m3 have been recorded at homes (Du et al., 2015; Xu et al., 2016; Zhong et al., 2017).


\paragraph{Formaldehyde}

φορμαλδεΰδη (HCHO) είναι ένα ιδιαίτερα δραστικό αέριο που εκλύεται κατά την ατελή αντίδραση υδρογονανθράκων. Η φορμαλδεΰδη παράγεται επίσης από την οξείδωση άλλων VOCs με όζον ή ακτινοβολία (Campagnolo et al., 2017; Kotzias et al., 2005; US Environmental Protection Agency, 1995; World Health Organization and WHO Regional Office for Europe, 2010). Αυτή η χημική ουσία χρησιμοποιείται ευρέως σε ρητίνες, κόλλες, χρώματα, προϊόντα χαρτιού, καλλυντικά, ηλεκτρονικό εξοπλισμό, καθαριστικά και υφάσματα. Η φορμαλδεΰδη είναι παρούσα σε δομικά υλικά, όπως οι μονωτικές αφρώδεις ύλες ξύλινων προϊόντων που χρησιμοποιούνται σε δάπεδα ή έπιπλα. Ωστόσο, οι εκπομπές από αυτά τα υλικά (π.χ. κόντρα πλακέ, μοριοσανίδα ή ίνες ξύλου) συνήθως μειώνονται μέσα σε μερικές εβδομάδες (Carazo Fernández et al., 2013; Jantunen et al., 1998; Kotzias et al., 2005; Royal College of Physicians, 2016; World Health Organization and WHO Regional Office for Europe, 2010).

Η φορμαλδεΰδη απορροφάται γρήγορα από το αναπνευστικό ή το γαστρεντερικό σύστημα. Οι οξείες επιδράσεις της έκθεσης περιλαμβάνουν οσμή, ερεθισμό, πονοκέφαλο και έκζεμα. Η φορμαλδεΰδη είναι γνωστό καρκινογόνο και γονοτοξικό χημικό, της οποίας η μακροχρόνια έκθεση μπορεί να προκαλέσει καρκίνο του ρινοφάρυγγα και μυελοειδή λευχαιμία (Kotzias et al., 2005; US Environmental Protection Agency, 1995; World Health Organization and WHO Regional Office for Europe, 2010).

Τα επίπεδα έκθεσης σε σπίτια, γραφεία και σχολεία είναι συγκρίσιμα (7,7–30, 8–17, 9–17 μg/m³) (Broderick et al., 2017; Canha et al., 2017; Dodson et al., 2008; Geiss et al., 2011; Kotzias et al., 2009; Langer et al., 2015; Mandin et al., 2017; Prasauskas et al., 2016; Zhong et al., 2017; Zielinska et al., 2012).

Formaldehyde (HCHO) is a highly reactive gas emitted during the incomplete reaction of hydrocarbons. Formaldehyde is also produced by oxidation of other VOCs with ozone or radiation (Campagnolo et al., 2017; Kotzias et al., 2005; US Environmental Protection Agency, 1995; World Health Organization and WHO Regional Office for Europe, 2010). This chemical is widely used in resins, glues, paints, paper products, cosmetics, electronic equipment, cleaning agents and fabrics. Formaldehyde is present in construction materials such as insulation foams of wooden-based materials employed in floorings or furniture. However, emissions from these materials (e.g. plywood, particleboard or fiberboard) usually decay within several weeks (Carazo Fernández et al., 2013; Jantunen et al., 1998; Kotzias et al., 2005; Royal College of Physicians, 2016; World Health Organization and WHO Regional Office for Europe, 2010). Formaldehyde is rapidly absorbed by the respiratory or gastrointestinal system. Acute exposure effects include odor, irritation, headache and eczema. Formaldehyde is a known carcinogenic and genotoxic chemical, whose long-term exposure can originate nasopharyngeal cancer and myeloid leukemia (Kotzias et al., 2005; US Environmental Protection Agency, 1995; World Health Organization and WHO Regional Office for Europe, 2010). The level of exposure in homes, offices and schools is comparable (7.7–30, 8–17, 9–17 μg/m3) (Broderick et al., 2017; Canha et al., 2017; Dodson et al., 2008; Geiss et al., 2011; Kotzias et al., 2009; Langer et al., 2015; Mandin et al., 2017; Prasauskas et al., 2016; Zhong et al., 2017; Zielinska et al., 2012).

\paragraph{Trichloroethylene}

Το τριχλωροαιθυλένιο (Trichloroethylene, TCE) είναι μια πτητική χημική ουσία που χρησιμοποιείται ευρέως ως διαλύτης και καθαριστικό μέσο (απολίπανση μετάλλων με ατμούς). Πηγές TCE σε εσωτερικούς χώρους περιλαμβάνουν προϊόντα στα οποία χρησιμοποιείται ως διαλύτης, όπως λιπαντικά, βερνίκια, αφαιρετικά χρωμάτων, κόλλες και υγρά διόρθωσης για γραφομηχανές. Το TCE μπορεί επίσης να υπάρχει σε ορισμένα οικιακά προϊόντα λεύκανσης και άλλα καθαριστικά. Επιπλέον, το TCE μπορεί να εντοπίζεται σε μολυσμένα ύδατα και εδάφη, τα οποία ενδέχεται να συμβάλλουν έμμεσα στα επίπεδα συγκέντρωσής του σε εσωτερικούς χώρους (Jantunen et al., 1998; Singh et al., 2016; World Health Organization and WHO Regional Office for Europe, 2010).

Η ανθρώπινη έκθεση σε αυτήν τη VOC πραγματοποιείται μέσω δερματικής απορρόφησης, κατάποσης και εισπνοής. Η έκθεση στο TCE είναι συχνά διαλείπουσα, καθώς τα προϊόντα που περιέχουν αυτή την ουσία δεν χρησιμοποιούνται συνεχώς. Η οξεία έκθεση στο TCE επηρεάζει το κεντρικό νευρικό σύστημα σε συγκεντρώσεις περίπου 270 mg/m³, προκαλώντας μείωση της αισθητηριακής ικανότητας. Το TCE είναι καρκινογόνος ουσία σε περίπτωση μακροχρόνιας έκθεσης, ενώ η χρόνια έκθεση μπορεί να προκαλέσει καρκίνο του ήπατος, των νεφρών και των χοληφόρων οδών, καθώς και λέμφωμα non-Hodgkin. Για τον λόγο αυτό, δεν συνιστάται ασφαλές όριο έκθεσης λόγω της καρκινογόνου δράσης του (World Health Organization and WHO Regional Office for Europe, 2010).

Συγκεντρώσεις TCE που κυμαίνονται από 0,3 έως 0,6 μg/m³ καταγράφονται συνήθως σε κατοικίες, ενώ σε εμπορικούς χώρους κυμαίνονται από 0,6 έως 4,7 μg/m³ (Dodson et al., 2008; Edwards et al., 2001; Mandin et al., 2017; Singh et al., 2016; Xu et al., 2016).

Trichloroethylene (TCE) is a volatile chemical widely used as solvent and cleaning agent (vapor degreasing of metals). Indoor sources include products where TCE is used as solvents such as lubricants, varnishes, paint removers, adhesives and typewriter correction fluids. TCE can be also present in some bleach household products and other cleaning agents. TCE can be present in contaminated waters and soils, which could indirectly contribute to indoor levels. (Jantunen et al., 1998; Singh et al., 2016; World Health Organization and WHO Regional Office for Europe, 2010). Human exposure to this VOC occurs by dermal absorption, ingestion and inhalation. Exposure to TCE is often intermittent, as the products that contain this pollutant are not constantly used. Acute exposure to TCE affects the central nervous system at concentrations of ∼270 mg/m3, causing a reduced sensorial capacity. TCE is a carcinogenic chemical under long-term exposure, and a chronic exposure may cause liver, kidney and bile duct cancer and non-Hodgkin’s lymphoma. Therefore, no safe threshold is recommended due to its carcinogenicity (World Health Organization and WHO Regional Office for Europe, 2010). TCE concentrations ranging from 0.3 to 0.6 μg/m3 are typically recorded at homes, and from 0.6 to 4.7 μg/m3 in shopping complexes (Dodson et al., 2008; Edwards et al., 2001; Mandin et al., 2017; Singh et al., 2016; Xu et al., 2016).

\paragraph{Alpha-pinene }

Ο α-πινένιο (άλφα-πινένιο) είναι ένα υγρό τερπένιο που απαντάται φυσικά στα φυτά και αποτελεί κοινό συστατικό των αιθέριων ελαίων. Το πινένιο χρησιμοποιείται ως διαλύτης σε ορισμένα χρώματα και υλικά στεγανοποίησης και μπορεί να βρεθεί σε προϊόντα αρωματοποιίας, αποσμητικά και καθαριστικά προϊόντα. Εκπομπές πινενίου μπορούν επίσης να προέρχονται από ξύλινα υλικά, όπως έπιπλα ή δάπεδα, ιδιαίτερα όταν αυτά είναι κατασκευασμένα από ξύλο πεύκου (Campagnolo et al., 2017; Kotzias et al., 2005; Luengas et al., 2015; Royal College of Physicians, 2016).

Τα διαθέσιμα δεδομένα σχετικά με τους κινδύνους για την υγεία από το πινένιο είναι περιορισμένα. Αν και η οξεία έκθεση σε υψηλές συγκεντρώσεις μπορεί να προκαλέσει ερεθισμό και φλεγμονή, δεν έχουν αναφερθεί μέχρι σήμερα κίνδυνοι από χρόνια έκθεση. Ωστόσο, ενδέχεται να προκύψουν κίνδυνοι για την υγεία όταν συνυπάρχει όζον ή άλλες δραστικές ρίζες, λόγω της υψηλής αντιδραστικότητας του πινενίου και άλλων ακόρεστων μορίων (Hubbard et al., 2005; Kotzias et al., 2005; Leung, 2015; SCHER, 2007). Τα παραπροϊόντα αυτών των αντιδράσεων μπορούν να προκαλέσουν ερεθισμό των οφθαλμών και των ανώτερων αεραγωγών σε μεγαλύτερο βαθμό και για μεγαλύτερη διάρκεια σε σύγκριση με το ίδιο το πινένιο.

Οι υψηλότερες συγκεντρώσεις πινενίου καταγράφονται σε κατοικίες (11–32 μg/m³) και βιβλιοθήκες (10–30 μg/m³), ενώ σε γραφεία και σχολεία παρατηρούνται συγκεντρώσεις έως 6,3 και 1,5 μg/m³ αντίστοιχα (Dodson et al., 2008; Edwards et al., 2001; Geiss et al., 2011; Langer et al., 2015; Mandin et al., 2017; Rösch et al., 2014; Xu et al., 2016; Zhong et al., 2017).

Alpha-pinene is a liquid terpene naturally present in plants and a common constituent of essential oils. Pinene is used as a solvent in some paints and waterproof substances, and can be found in perfumery products, deodorizers and cleaning products. Pinene emissions can also originate from wooden materials like furniture or flooring, particularly those made from pine wood (Campagnolo et al., 2017; Kotzias et al., 2005; Luengas et al., 2015; Royal College of Physicians, 2016). Limited data of pinene health risks is available. While acute exposure to high concentration can produce irritation and inflammation, no chronic exposure risks have been reported yet. Health risks might be expected when ozone or other reactive radicals are present, due to their high reactivity with pinene and other unsaturated molecules (Hubbard et al., 2005; Kotzias et al., 2005; Leung, 2015; SCHER, 2007). These reaction by-products would produce eye and upper airways irritation in a more extended way than pinene itself. The largest pinene concentrations are found at homes (11–32 μg/m3) and libraries (10–30 μg/m3), while offices and schools exhibit concentrations up to 6.3 and 1.5 μg/m3, respectively (Dodson et al., 2008; Edwards et al., 2001; Geiss et al., 2011; Langer et al., 2015; Mandin et al., 2017; Rösch et al., 2014; Xu et al., 2016; Zhong et al., 2017).

\paragraph{Limonene}

Το λιμονένιο είναι ένα φυσικό τερπένιο που υπάρχει σε δύο ισομερείς μορφές (d- και l-). Χρησιμοποιείται σε μια μεγάλη ποικιλία οικιακών προϊόντων, όπως καθαριστικά, ρητίνες, αρωματικά χώρου, αποσμητικά, αρώματα και σαμπουάν (Campagnolo et al., 2017; Kotzias et al., 2005; Luengas et al., 2015; Royal College of Physicians, 2016). Το τερπένιο αυτό βρίσκεται επίσης στα τρόφιμα ως πρόσθετο λόγω της κιτρικής του οσμής και γεύσης.

Η έκθεση στο λιμονένιο γίνεται κυρίως μέσω εισπνοής ή κατάποσης. Ο συγκεκριμένος ρύπος έχει χαμηλή οξεία τοξικότητα και μέχρι στιγμής δεν έχουν περιγραφεί συμπτώματα πέρα από ερεθισμό των ματιών ή του δέρματος. Δεν έχουν μελετηθεί οι μακροπρόθεσμες επιπτώσεις στην υγεία από το λιμονένιο και δεν υπάρχει ένδειξη καρκινογένεσης ή γονοτοξικότητας. Οι κίνδυνοι για την υγεία συνδέονται κυρίως με την αντίδραση του λιμονενίου με το όζον ή άλλες δραστικές ρίζες στο περιβάλλον, που δημιουργούν παραπροϊόντα όπως αλδεΰδες, καρβοξυλικά οξέα και υπεροξείδια, τα οποία ευθύνονται για ερεθισμό και δυσάρεστη οσμή ακόμα και σε χαμηλές συγκεντρώσεις (Hubbard et al., 2005; Kotzias et al., 2005; Leung, 2015; SCHER, 2007).

Limonene is a natural terpene present as two isomers (d- and l-). Limonene is used in a wide variety of household products like cleaning agents, resins, air fresheners, deodorants, fragrances and shampoos (Campagnolo et al., 2017; Kotzias et al., 2005; Luengas et al., 2015; Royal College of Physicians, 2016). This terpene is also found in food as an additive due to its citric odor and flavor. Limonene intake occurs by inhalation or ingestion. This pollutant has low acute toxicity, no symptoms further than eye or skin irritation have been described. Chronic health effects for limonene have not been studied, and there is no evidence of carcinogenicity or genotoxicity. Heath risks are associated to ambient reaction with ozone or reactive radicals, which form byproducts such as aldehydes, carboxylic acids and peroxides, responsible of irritation and odor annoyance at low concentrations (Hubbard et al., 2005; Kotzias et al., 2005; Leung, 2015; SCHER, 2007). Maximum concentrations of 32, 19 and 11 μg/m3 have been recorded at homes, offices and schools, respectively (Dodson et al., 2008; Du et al., 2015; Edwards et al., 2001; Geiss et al., 2011; Langer et al., 2015; Mandin et al., 2017; Rösch et al., 2014; Xu et al., 2016; Zhong et al., 2017).



\cite{GONZALEZMARTIN2021128376}


\subsubsection{Διοξείδιο του αζώτου}

Το μονοξείδιο του αζώτου (NO) και το διοξείδιο του αζώτου (NO₂) είναι οι πιο συνηθισμένες και σημαντικές ενώσεις οξειδίων του αζώτου. Το NO παράγεται σε μεγαλύτερο βαθμό και οξειδώνεται περαιτέρω σε NO₂ παρουσία οξυγόνου, όζοντος ή VOCs (European Environment Agency, 2017; World Health Organization and WHO Regional Office for Europe, 2010).

Οι πιο συχνοί εσωτερικοί (indoor) ρύποι NOx προέρχονται από συσκευές αερίου, όπως κουζίνες, φούρνους ή θερμοσίφωνες. Το ETS (παθητικό κάπνισμα) και τα τζάκια μπορούν επίσης να συμβάλλουν στις εκπομπές NOx (Carazo Fernández et al., 2013; Kotzias et al., 2005; Leung, 2015; Luengas et al., 2015; Royal College of Physicians, 2016; Singh et al., 2016; US Environmental Protection Agency, 1995; World Health Organization and WHO Regional Office for Europe, 2010). Η εισχώρηση αέρα από το εξωτερικό επηρεάζει σημαντικά τα εσωτερικά επίπεδα, ιδιαίτερα κοντά σε δρόμους ή περιοχές υψηλής βιομηχανικής πυκνότητας.

Το εισπνεόμενο NO₂ μετατρέπεται σε νιτρικό οξύ στους πνεύμονες, προκαλώντας βλάβες στα κύτταρα και στο ανοσοποιητικό σύστημα. Η εισπνοή μπορεί να προκαλέσει αναπνευστική δυσλειτουργία όταν υπάρχουν προϋπάρχουσες παθήσεις (όριο προστασίας για άτομα με άσθμα: 200 μg/m³ μέσος όρος 1 ώρας, ≈2000 μg/m³ για φυσιολογικά άτομα). Το NO₂ μπορεί επίσης να προκαλέσει ερεθισμό των ματιών από άμεση επαφή (Carazo Fernández et al., 2013; Kotzias et al., 2005; Leung, 2015; US Environmental Protection Agency, 1995; World Health Organization and WHO Regional Office for Europe, 2010).

Συγκεντρώσεις 5–17 μg/m³, 16–18 μg/m³ και έως 30 μg/m³ έχουν μετρηθεί σε σπίτια, γραφεία και εμπορικές περιοχές, αντίστοιχα (Broderick et al., 2017; Langer et al., 2015; Mandin et al., 2017; Prasauskas et al., 2016; Singh et al., 2016; Vasile et al., 2016).

Nitrogen monoxide (NO) and nitrogen dioxide (NO2) are the most common and important nitrogen oxides. NO is produced in a larger extent and is further oxidized to NO2 in the presence of oxygen, ozone or VOCs (European Environment Agency, 2017; World Health Organization and WHO Regional Office for Europe, 2010). The most frequent indoor sources of NOx are gas appliances like stoves, ovens or water heaters. ETS and fireplaces can also contribute to NOx emissions (Carazo Fernández et al., 2013; Kotzias et al., 2005; Leung, 2015; Luengas et al., 2015; Royal College of Physicians, 2016; Singh et al., 2016; US Environmental Protection Agency, 1995; World Health Organization and WHO Regional Office for Europe, 2010). Infiltration from outdoor air strongly influences indoor levels, in particular within short distance from roadways or high-density industrial areas. Inhaled NO2 is converted to nitric acid in the lungs, with the associate damage to cells and to the immune system. Inhalation can cause respiratory malfunctioning when other pathologies are previously present (limit for protection of asthmatic of 200 μg/m3 1-h average; ≈2000 μg/m3 impact normal individuals). NO2 can cause also eye irritation by direct contact (Carazo Fernández et al., 2013; Kotzias et al., 2005; Leung, 2015; US Environmental Protection Agency, 1995; World Health Organization and WHO Regional Office for Europe, 2010). Concentrations of 5–17 μg/m3, 16–18 μg/m3 and up to 30 μg/m3 have been measured in houses, offices and shopping areas, respectively (Broderick et al., 2017; Langer et al., 2015; Mandin et al., 2017; Prasauskas et al., 2016; Singh et al., 2016; Vasile et al., 2016).

\cite{GONZALEZMARTIN2021128376}


.2. Διοξείδιο του Αζώτου (NO₂)
3.2.1. Επίπεδα Έκθεσης, Πηγές και Παράγοντες Προσδιορισμού

Το NO₂ αποτελεί προϊόν καύσης που παράγεται από οχήματα, παραγωγή ενέργειας και άλλες εξωτερικές πηγές καύσης, καθώς και από εσωτερικές πηγές όπως συσκευές αερίου και θερμαντήρες κεροζίνης (π.χ., [55,56]).

Τα επίπεδα εσωτερικού NO₂ που αναφέρθηκαν σε σαράντα έξι μελέτες κυμαίνονταν μεταξύ 3,4 μg/m³ σε σπίτια στην πόλη Quebec, Καναδάς [27], και 1210 μg/m³ σε σπίτια με εσωτερικό κάπνισμα και/ή μαγείρεμα στο Λάο, PDR [38], ξεπερνώντας σε ορισμένες περιπτώσεις τις οδηγίες του ΠΟΥ [1] για οξεία (200 μg/m³) ή χρόνια (40 μg/m³) έκθεση (Πίνακας S1). Τα επίπεδα εσωτερικού NO₂ όπου υπήρχαν πηγές όπως συσκευές αερίου ήταν σημαντικά υψηλότερα σε σχέση με τα εξωτερικά, όπου η κύρια πηγή ήταν η κυκλοφορία οχημάτων [57].

Οι συγκεντρώσεις NO₂ ήταν υψηλότερες σε όλα τα δωμάτια των σπιτιών στο Ηνωμένο Βασίλειο [58] και στη Γερμανία [59], όπου χρησιμοποιούνταν κουζίνες αερίου. Τα επίπεδα NO₂ στις κουζίνες με κουζίνα αερίου στο Ηνωμένο Βασίλειο ήταν διπλάσια σε σχέση με κουζίνες με ηλεκτρική κατά τη διάρκεια του χειμώνα [58]. Τα οξείδια του αζώτου (NOx) και το NO₂ ήταν υψηλότερα σε σπίτια με πιλότους στις εστίες αερίου σε σύγκριση με μαγείρεμα με αέριο χωρίς πιλότο. Σε σπίτια στην Καλιφόρνια, ΗΠΑ, όπου οι κάτοικοι μαγείρευαν 4 ώρες ή περισσότερες την ημέρα με αέριο, η χρήση απορροφητήρων κουζίνας συνδεόταν με χαμηλότερη έκθεση σε NOx και NO₂ [60].

Επιπλέον, παρατηρήθηκαν εποχιακές επιδράσεις με αυξημένες συγκεντρώσεις NO₂ το χειμώνα στην Ιαπωνία [55], πιθανώς λόγω χρήσης θερμαντήρων αερίου. Οι Cyrys et al. [59] παρατήρησαν εξωτερικές συγκεντρώσεις NO₂ περίπου διπλάσιες των εσωτερικών στη Γερμανία, ενώ οι Cibella et al. [61] και Zipprich et al. [62] ανέφεραν σημαντικά υψηλότερα επίπεδα NO₂ σε εσωτερικούς χώρους σε Ιταλία και ΗΠΑ αντίστοιχα. Συνολικά, οι πιο σημαντικοί παράγοντες πρόβλεψης των εσωτερικών συγκεντρώσεων NO₂ ήταν η χρήση κουζίνας αερίου, ακολουθούμενη από τον εξαερισμό και τα επίπεδα NO₂ εξωτερικού αέρα.

3.2.2. Επιπτώσεις στην Υγεία

Τα περισσότερα στοιχεία, κυρίως από μελέτες στις ΗΠΑ, υποδεικνύουν ότι η έκθεση σε υψηλότερες συγκεντρώσεις εσωτερικού NO₂ προκαλεί συμπτώματα σε παιδιά με άσθμα, όπως σφίξιμο στο στήθος, δύσπνοια, συριγμό, βήχα, νυχτερινά συμπτώματα, αυξημένο αριθμό κρίσεων άσθματος και χρήση εισπνευστήρων, καθώς και μειωμένο εκπνεόμενο όγκο σε ένα δευτερόλεπτο (FEV₁) [63,64,65,66,67].

Τα παιδιά που ζουν σε αστικά κέντρα φαίνεται να διατρέχουν υψηλότερο κίνδυνο για τις αρνητικές επιπτώσεις του NO₂ λόγω της σχετικά υψηλής εσωτερικής έκθεσής τους [64], αν και αυξημένος κίνδυνος νοσηρότητας από άσθμα παρατηρήθηκε και σε συγκεντρώσεις NO₂ κοινές σε αστικά και προαστιακά σπίτια [57]. Επιπλέον, η έκθεση σε υψηλότερες συγκεντρώσεις εσωτερικού NO₂ συνδέθηκε με αυξημένα αναπνευστικά συμπτώματα και κίνδυνο έξαρσης ΧΑΠ σε πρώην καπνιστές με μέτρια έως σοβαρή ΧΑΠ [56].


 3.2. Nitrogen Dioxide
3.2.1. Exposure Levels, Sources and Determinants
NO2 is a by-product of combustion produced by motor vehicles, energy generation and other outdoor sources involving combustion, as well as indoor sources such as gas appliances and kerosene heaters (e.g., [55,56]).
Indoor NO2 levels reported in forty-six studies ranged between 3.4 μg/m3 in homes in Quebec City, Canada [27], and 1210 μg/m3 in houses with indoor smoking and/or cooking in Lao PDR [38], in cases exceeding the WHO guidelines [1] for acute (200 μg/m3) or chronic (40 μg/m3) exposure (Table S1). Indoor levels where NO2 sources, such as gas appliances, were present were much higher than outdoors, where the primary source of NO2 was road traffic [57]. NO2 concentrations were higher in all rooms in houses in the UK [58] and Germany [59] where gas cookers were used. NO2 levels in kitchens with a gas cooker in the UK were twice as high as in those with an electric cooker during winter [58]. Nitrogen oxides (NOx) and NO2 were higher in homes with cooktop pilot burners, relative to gas cooking without pilots. In homes in California, USA, where residents cooked 4 h or more per day with gas, self-reported use of kitchen exhaust fans was associated with lower NOx and NO2 exposure [60].
Seasonal effects were also identified with increased indoor concentrations of NO2 in winter in Japan [55], which was thought to be due to the use of gas heaters. Cyrys et al. [59] observed outdoor NO2 concentrations approximately twice as high as indoor levels in Germany, while Cibella et al. [61] and Zipprich et al. [62] observed significantly higher indoor NO2 concentrations than outdoor concentrations in Italy and USA respectively. Overall, the most important predictors of indoor NO2 concentrations were gas cooker use, followed by ventilation and outdoor NO2 levels.
3.2.2. Health Effects
Most reviewed evidence, mainly from US studies, suggested that exposure to higher indoor NO2 concentrations leads to symptoms in children with asthma, including chest tightness, shortness of breath, wheeze, cough, nocturnal symptoms, increased number of asthma attacks and inhaler use, and decreased forced expiratory volume in one second (FEV1) [63,64,65,66,67]. Children living in inner cities appeared to be at higher risk for the adverse effects of NO2 given their relatively high indoor exposure [64], although increased risk of asthma morbidity also occurred at NO2 concentrations common in urban and suburban homes [57]. In addition, exposure to higher indoor NO2 concentrations was associated with increased respiratory symptoms and risk of COPD exacerbations in former smokers with moderate to severe COPD [56].

\cite{ijerph17238972}

3.3. Οξείδια του Αζώτου (NOx)

Τα δύο κύρια οξείδια του αζώτου είναι το μονοξείδιο του αζώτου (NO) και το διοξείδιο του αζώτου (NO₂), τα οποία σχετίζονται κυρίως με πηγές καύσης, όπως οι κουζίνες και οι θερμαντήρες [46]. Οι συγκεντρώσεις NO και NO₂ στο περιβάλλον ποικίλλουν σημαντικά ανάλογα με τις τοπικές πηγές και καταβόθρες.

Η μέση συγκέντρωσή τους σε κτίρια χωρίς δραστηριότητες καύσης είναι περίπου η μισή από εκείνη του εξωτερικού αέρα, αλλά όταν χρησιμοποιούνται κουζίνες και θερμαντήρες αερίου, τα εσωτερικά επίπεδα συχνά υπερβαίνουν τα εξωτερικά. Σε περιβαλλοντικές συνθήκες, το NO οξειδώνεται γρήγορα για να σχηματίσει NO₂· ως εκ τούτου, το NO₂ θεωρείται συνήθως ως πρωτογενής ρύπος.

Η αντίδραση του NO₂ με το νερό παράγει νιτρώδες οξύ (HONO), έναν ισχυρό οξειδωτικό παράγοντα και κοινό ρύπο σε εσωτερικά περιβάλλοντα [47]. Έχει αναφερθεί ότι τα εσωτερικά επίπεδα NO₂ εξαρτώνται τόσο από εξωτερικές όσο και από εσωτερικές πηγές· συνεπώς, τα εσωτερικά επίπεδα μπορούν να επηρεαστούν από υψηλές εξωτερικές συγκεντρώσεις που προέρχονται από καύση ή τοπικές πηγές κυκλοφορίας. Επίσης, αναφέρθηκε ότι η απόσταση μεταξύ κτιρίων και δρόμων έχει σημαντική επίδραση στα εσωτερικά επίπεδα NO₂ [48].

Επιπλέον, η ανταλλαγή αέρα μεταξύ εξωτερικού και εσωτερικού χώρου επηρεάζει επίσης τα επίπεδα NO₂ στα κτίρια [49]. Οι κύριες εσωτερικές πηγές περιλαμβάνουν το κάπνισμα και συσκευές καύσης ξύλου, αερίου, πετρελαίου, άνθρακα και κεροζίνης, όπως κουζίνες, θερμαντικά σώματα, φούρνοι, θερμοσίφωνες και τζάκια [47].

 3.3. NOX
The two principal nitrogen oxides are nitric oxide (NO) and nitrogen dioxide (NO2), both of which are associated with combustion sources, such as cooking stoves and heaters [46]. Ambient concentrations of NO and NO2 vary widely depending on local sources and sinks. Their average concentration in buildings without combustion activities is half that in the outdoors, but when gas stoves and heaters are used, indoor levels often exceed outdoor levels. Under ambient conditions, NO is rapidly oxidized to form NO2; hence, NO2 is usually considered as a primary pollutant. The reaction of NO2 with water produces nitrous acid (HONO), a strong oxidant and common pollutant of indoor environments [47]. It has been indicated that indoor levels of NO2 are a function of both outdoor and indoor sources; hence, indoor levels can be influenced by high outdoor levels originating from combustion or local traffic sources. It was reported that the distance between buildings and roadways has a significant influence on indoor NO2 levels [48]. Besides, the air exchange between outdoors and indoors also affects NO2 levels in buildings [49]. Additionally, the major indoor sources include smoking and wood-, gas-, oil-, coal-, and kerosene-burning appliances, such as stoves, space, ovens, and water heaters and fireplaces [47].



\cite{ijerph17238972}

NO2 is a highly reactive gas which is related to the development of ozone and PM2.5. NO2 primarily gets into the air from the burning of fuel. Similarly to sulfur dioxide, it can cause respiratory symptoms and airway inflammation. Asthmatics, children, and older adults are at higher risk from this pollutant [60].

\cite{su12219045}

\subsubsection{CO2 KAI C0}

 3.6. COx
Carbon monoxide (CO) in indoor air is produced mainly by combustion processes, such as cooking or heating. Besides, CO can also enter into indoor environments through infiltration from outdoor air [63]. The important sources of indoor CO emissions include unvented kerosene and gas space heaters; leaking chimneys and furnaces; back-drafting from furnaces, gas water heaters, wood stoves, and fireplaces; gas stoves; generators and other gasoline-powered equipment; and tobacco smoke [47]. The average concentration of CO in a building without any gas stoves is about 0.5–5 ppm, while the concentration in areas near gas stoves ranges from 5 to 15 ppm and even 30 ppm or higher. CO exposure can cause adverse health effects, such as (i) at low concentrations, there are impacts on cardiovascular and neurobehavioral processes; and (ii) at high concentrations, unconsciousness or death [64].
Carbon dioxide (CO2), a colorless and odorless gas, is a well-known constituent of the earth’s atmosphere and also a major human metabolite [65]. The average CO2 concentration in ambient air is about 400 ppm, which is primarily the result of the combustion of fossil fuels [65,66]. Recently, the indoor CO2 level has been applied as a reference for the assessment of IAQ as well as for ventilation control [66,67,68]. According to the ASHRAE standard, it is recommended that indoor CO2 concentrations are below 700 ppm to ensure human health [69]. It is established that exposure to a a CO2 concentration of 3000 ppm increases headache intensity, sleepiness, fatigue, and concentration difficulty [65,70].

\cite{ijerph17238972}

Carbon monoxide is a toxic gas emitted as a result of incomplete combustion processes. The main indoor sources of CO are tobacco smoke, defective cooking and heating devices, fireplaces and vehicle gases from attached garages, also including outdoor air exchange in dense traffic areas or highly industrialized districts. (Carazo Fernández et al., 2013; Kotzias et al., 2005; Leung, 2015; Royal College of Physicians, 2016; Singh et al., 2016; US Environmental Protection Agency, 1995; World Health Organization and WHO Regional Office for Europe, 2010). Inhalation is the main entrance to human body of CO, which is able to bind reversibly to hemoglobin with an affinity 200–250 times higher than oxygen. CO toxic effects are thus based on the interference with O2 transportation through the human body resulting in tissue hypoxia. The extent of health damage depends on the concentration and duration of exposure, and include cardiovascular, respiratory and neurological problems. CO can even cause death under high acute exposure (>1000 ppmv) or long-term exposure (Kotzias et al., 2005; Leung, 2015; US Environmental Protection Agency, 1995; World Health Organization and WHO Regional Office for Europe, 2010). Concentrations of 0.5–17 ppmv, 0.9–3.8 ppmv and up to 8.5 ppmv have been recorded in kitchens, bedrooms and shopping areas, respectively (Broderick et al., 2017; Canha et al., 2017; Sidhu et al., 2017; Singh et al., 2016; Vasile et al., 2016; Yamamoto et al., 2014; Yip et al., 2017; Zielinska et al., 2012).


\cite{GONZALEZMARTIN2021128376}

 \subsubsection{3.7. Toxic Metals}

 Βαρέα Μέταλλα

Τα βαρέα μέταλλα απελευθερώνονται στην ατμόσφαιρα είτε μέσω ανθρώπινων δραστηριοτήτων είτε μέσω φυσικών διεργασιών [71]. Η εσωτερική ατμοσφαιρική ρύπανση (IAP) από βαρέα μέταλλα έχει πολλαπλές αιτίες, όπως διείσδυση εξωτερικών ρύπων (σκόνη και έδαφος), κάπνισμα, προϊόντα κατανάλωσης καυσίμων και οικοδομικά υλικά [55].

Τα βαρέα μέταλλα που περιέχονται στη σκόνη εσωτερικών χώρων μπορούν να εισέλθουν στον ανθρώπινο οργανισμό μέσω εισπνοής, κατάποσης ή επαφής με το δέρμα και να έχουν επιβλαβείς επιπτώσεις στην υγεία [72,73]. Σύμφωνα με τον Διεθνή Οργανισμό Έρευνας για τον Καρκίνο (IARC), τα βαρέα μέταλλα στον εσωτερικό αέρα ταξινομούνται σε δύο κύριες ομάδες ανάλογα με τις επιδράσεις τους στον άνθρωπο:

Μη καρκινογόνα στοιχεία, όπως κοβάλτιο (Co), αλουμίνιο (Al), χαλκός (Cu), νικέλιο (Ni), σίδηρος (Fe) και ψευδάργυρος (Zn).

Στοιχεία που είναι είτε καρκινογόνα είτε μη καρκινογόνα, όπως αρσένιο (As), χρώμιο (Cr), κάδμιο (Cd) και μόλυβδος (Pb) [74].

Τα πιο κοινά βαρέα μέταλλα (As, Cr, Cd, Pb) μπορούν να προκαλέσουν καρκίνο [75,76], ενώ το Cd και το Pb, μαζί με κάποια άλλα, μπορεί να προκαλέσουν και άλλες βλαβερές επιδράσεις, όπως καρδιαγγειακές παθήσεις, καθυστέρηση ανάπτυξης και βλάβες στο νευρικό σύστημα [73,77,78].

Έχει αναφερθεί ότι τα επίπεδα Pb στον εσωτερικό αέρα κυμαίνονται από 5,80 έως 639,10 μg/g, ενώ τα μέγιστα επίπεδα των As, Al, Cr, Cd, Co, Cu, Ni, Fe και Zn είναι περίπου 486,80, 7150,00, 254,00, 8,48, 43,40, 513,00, 471,00, 4801,00 και 2293,56 μg/g αντίστοιχα [74].

\cite{ijerph17238972}



Heavy metals are released into the atmosphere through either human activities or natural processes [71]. IAP by heavy metals has various causes, including infiltration of outdoor pollutants (dust and soil), smoking, fuel consumption products, and building materials [55]. Heavy metals in indoor dust, entering the human body through inhalation, ingestion, or dermal contact, can have adverse effects on human health [72,73]. According to the International Agency for Research on Cancer (IARC), heavy metals in indoor air are classified into two main groups based on their effects on humans: (i) Non-carcinogenic elements, including cobalt (Co), aluminum (Al), copper (Cu), nickel (Ni), iron (Fe), and zinc (Zn); and (ii) both carcinogenic and non-carcinogenic elements encompassing arsenic (As), chromium (Cr), cadmium (Cd), and lead (Pb) [74]. These common heavy metals (i.e., As, Cr, Cd, Pb) are likely to cause cancers [75,76], while Cd and Pb, along with some others, can cause carcinogenic effects, such as cardiovascular disease, slow growth development, and damage to the nervous system [73,77,78]. It has been reported that Pb levels in indoor air can fluctuate from 5.80 to 639.10 μg/g, while the highest levels of As, Al, Cr, Cd, Co, Cu, Ni, Fe, and Zn are about 486.80, 7150.00, 254.00, 8.48, 43.40, 513.00, 471.00, 4801.00, and 2293.56 μg/g, respectively [74]

\cite{ijerph17238972}

\subsubsection{PAH}

3.5 Πολυκυκλικοί Αρωματικοί Υδρογονάνθρακες (PAHs)
Έκθεση, Πηγές και Παράγοντες

Οι PAHs είναι προϊόντα ατελούς καύσης που παράγονται μέσω της καύσης ξύλου, άνθρακα, πετρελαίου και αερίου, του καπνίσματος, της καύσης αποβλήτων, της βιομηχανικής παραγωγής ενέργειας και των εκπομπών οχημάτων [75]. Οι PAHs χαμηλού μοριακού βάρους (δηλαδή με 2 ή 3 δακτυλίους) εκπέμπονται στη αέρια φάση, ενώ οι PAHs υψηλού μοριακού βάρους (με 5 ή περισσότερους δακτυλίους) εκπέμπονται κυρίως στη σωματιδιακή φάση [109].

Δώδεκα μελέτες ανέφεραν επίπεδα PAHs σε εσωτερικούς χώρους, με το άθροισμα διαφορετικών ομάδων PAHs να κυμαίνεται από 1,5 ng/m³ σε σπίτια στο Agra της Ινδίας [40], έως 9568 ng/m³ σε σπίτια στο Hangzhou της Κίνας [110] (Πίνακας S1). Πολλές μελέτες, κυρίως από Κίνα και Ιαπωνία, έδειξαν ότι οι συγκεντρώσεις PAHs στους εσωτερικούς χώρους των κατοικιών είναι γενικά υψηλότερες από τις εξωτερικές [110,111], αν και το αντίθετο παρατηρήθηκε σε ένα προάστιο στο Brisbane της Αυστραλίας [112]. Οι PAHs στον εσωτερικό αέρα, ιδιαίτερα οι χαμηλού μοριακού βάρους, προέρχονται κυρίως από εσωτερικές πηγές εκπομπής. Υψηλότερες συγκεντρώσεις PAHs παρατηρήθηκαν σε σπίτια με καπνιστές στην Κίνα, την Ιαπωνία και το Ηνωμένο Βασίλειο [84,113].

Σε σπίτια χωρίς καπνιστές, οι απωθητικές ουσίες για έντομα και οι πρακτικές μαγειρέματος ήταν οι κύριες πηγές των 2- και 3-δακτυλίων PAHs αντίστοιχα. Οι PAHs χαμηλού μοριακού βάρους συνδέονταν επίσης με τη χρήση καυστήρων κηροζίνης και τη ρύπανση εξωτερικού αέρα. Σε μια μελέτη στην Ιταλία, η παρουσία τζακιών εξηγούσε τα υψηλότερα επίπεδα συνολικών PAHs και πυρίνης (4-δακτυλικοί PAHs) σε αγροτικές κατοικίες σε σχέση με αστικές [114]. Οι PAHs υψηλού μοριακού βάρους (5- και 6-δακτυλικοί) στον εσωτερικό αέρα συνδέονταν κυρίως με εξωτερικές πηγές και τα επίπεδά τους ήταν ίσα ή χαμηλότερα από αυτά του εξωτερικού αέρα.

Στην Ινδία, οι συγκεντρώσεις PAHs τόσο στον εσωτερικό όσο και στον εξωτερικό αέρα ήταν σημαντικά υψηλότερες στην αέρια φάση (2-, 3-, 4-δακτυλικοί PAHs) σε σχέση με τη σωματιδιακή φάση (5- και 6-δακτυλικοί PAHs) (Masih et al., 2010). Στους εσωτερικούς χώρους, οι PAHs οφείλονται κυρίως στη χρήση αερίου, το μαγείρεμα (τηγάνισμα και καύση λαδιού), το κάπνισμα και το κάψιμο λιβανιού, ενώ στους εξωτερικούς χώρους οι πιο κοινές πηγές είναι τα βενζινοκίνητα και ντιζελοκίνητα οχήματα.

Επιπτώσεις στην Υγεία

Υπάρχουν αποδείξεις για τη γονοτοξικότητα και την καρκινογένεση πολλών PAHs σε ζωικά είδη, και επιδημιολογικές μελέτες έδειξαν συσχέτιση μεταξύ της έκθεσης σε PAHs και της εμφάνισης καρκίνου σε διάφορους ανθρώπινους ιστούς [115].

\cite{ijerph17238972}

 3.5. Polycyclic Aromatic Hydrocarbons
3.5.1. Exposure Levels, Sources and Determinants
PAHs are products of incomplete combustion generated through wood, coal, oil, and gas burning, smoking, waste incineration, industrial power generation, and vehicle emissions (e.g., [75]). PAHs with low molecular weight (i.e., with 2 or 3 rings) are emitted in the gaseous phase, while high-molecular-weight PAHs (i.e., with 5 or more rings) are emitted in the particulate phase [109].
Twelve studies reported indoor levels of PAHs, with the sum of different groups of PAHs ranging from 1.5 ng/m3 in homes in Agra, India [40], to 9568 ng/m3 in homes in Hangzhou, China [110] (Table S1). A number of studies, mostly from China and Japan, have shown residential indoor PAH concentrations generally higher than outdoor [110,111], but the opposite was observed in a suburban home in Brisbane, Australia [112]. Indoor air PAHs, especially low-molecular-weight PAHs, mainly came from indoor emission sources. Higher PAH concentrations were observed in smoking houses in China, Japan, and the UK [84,113]. In non-smoking houses, moth repellent and cooking practice were the main sources of 2- and 3-ring PAHs, respectively. Low-molecular-weight PAHs were also associated with kerosene heating and the outdoor pollution in non-smoking houses. In a study in Italy, the presence of fireplaces could explain the higher indoor levels of total PAHs and pyrene (4-ring PAHs) found in rural households compared to urban households [114]. High-molecular-weight (5- and 6-ring) PAHs in indoor air were mainly associated with outdoor sources and their levels tended to be the same or lower than those in outdoor air.
PAH concentrations both indoor and outdoor were significantly higher in the gaseous fraction (2-, 3-, and 4-ring PAHs) than in the particulate fraction (5- and 6-ring PAHs) in India (Masih et al., 2010). PAHs in the indoor environment were mainly attributable to gas usage, cooking (frying and oil combustion), smoking and incense burning, whereas outdoors the most common sources of PAHs were petrol and diesel vehicles.
3.5.2. Health Effects
There is evidence for the genotoxicity and carcinogenicity of many PAHs in animal species, and epidemiological studies demonstrated that there is a correlation between PAH exposure and cancer incidence for various human tissues [115].

\cite{ijerph17238972}


\subsubsection{Ozone}

Το όζον είναι ένας ισχυρός οξειδωτικός παράγοντας, που παράγεται κυρίως μέσω φωτοχημικών αντιδράσεων του O₂, των NOx και των VOCs στην ατμόσφαιρα. Ωστόσο, δεν μπορεί να χρησιμοποιηθεί για την εξουδετέρωση άλλων χημικών ρύπων στους εσωτερικούς χώρους, λόγω της αργής αντίδρασής του με τους περισσότερους ατμοσφαιρικούς ρύπους [50,51].

Το όζον μπορεί να αντιδράσει γρήγορα με αρκετούς εσωτερικούς ρύπους, αλλά τα προϊόντα των αντιδράσεων αυτών μπορούν να προκαλέσουν ερεθισμούς στον άνθρωπο και φθορά στα υλικά. Οι κύριες πηγές εσωτερικού όζοντος προέρχονται κυρίως από:

την εξωτερική ατμόσφαιρα,τη λειτουργία ηλεκτρικών συσκευών [52].Οι συσκευές που συχνά εκπέμπουν όζον στους εσωτερικούς χώρους περιλαμβάνουν φωτοαντιγραφικά μηχανήματα, συσκευές απολύμανσης, καθαρισμού αέρα και άλλες γραφειακές συσκευές [53,54,55,56]. Οι μηχανισμοί εκπομπής όζοντος από αυτές τις συσκευές χωρίζονται σε δύο κατηγορίες:
Εκφόρτιση κορώνης (Corona discharge) Φωτοχημικοί μηχανισμοί
Έχει δείξει ότι τα επίπεδα του όζοντος στους εσωτερικούς χώρους εξαρτώνται από διάφορους παράγοντες:

Τα επίπεδα όζοντος στο εξωτερικό περιβάλλον,Τα ποσοστά εκπομπής στο εσωτερικό,Τα ποσοστά ανταλλαγής αέρα.Τα ποσοστά απομάκρυνσης από επιφάνειες,Τις αντιδράσεις με άλλες χημικές ουσίες στον αέρα [50]

Συνήθως, τα επίπεδα όζοντος στους εσωτερικούς χώρους κυμαίνονται μεταξύ 20\% και 80\% των επιπέδων του εξωτερικού όζοντος, ανάλογα με το ποσοστό ανταλλαγής αέρα [57]. Ο άνθρωπος εκτίθεται κυρίως μέσω της εισπνοής, αλλά η επαφή με το δέρμα αναγνωρίζεται επίσης ως διαδρομή έκθεσης [58].

Ozone is a powerful oxidizing agent mainly produced by photochemical reactions of O2, NOx, and VOCs in the atmosphere. However, it cannot be used to eliminate other indoor chemical pollutants, due to its slow reaction with most airborne pollutants [50,51]. Ozone enables rapid reaction with several indoor pollutants, but the reaction products can irritate humans and damage materials. The main sources of indoor ozone mainly come from the outdoor atmosphere and the operation of electrical devices [52]. The machines commonly emitting indoor ozone gas include photocopiers, disinfecting devices, air-purifying devices, and other office devices [53,54,55,56]. The ozone emission mechanisms of these devices can be divided into two categories: Corona discharge and photochemical mechanisms. It has been shown that indoor ozone levels depend on various factors: (i) The outdoor ozone level; (ii) indoor emission rates; (iii) air-exchange rates; (iv) surface-removal rates, and (v) reactions between other chemicals and ozone in the air [50]. Indoor ozone levels generally fluctuate between 20% and 80% of the outdoor ozone level according to the air-exchange rate [57]. Humans are exposed to ozone primarily by inhalation, but skin exposure is also a recognized vector [58]. 
\cite{ijerph17238972}

\subsubsection{SO2}

Το διοξείδιο του θείου (SO₂) είναι το πιο κοινό αέριο στην ομάδα των θειωδών οξειδίων (SOx) που βρίσκονται στην ατμόσφαιρα. Το SO₂ παράγεται κυρίως από διαδικασίες καύσης ορυκτών καυσίμων και συνδέεται με αερολύματα και σωματίδια (PM) για να σχηματίσει ένα πολύπλοκο σύνολο ατμοσφαιρικών ρύπων [59].

Οι εσωτερικές πηγές εκπομπής SO₂ περιλαμβάνουν θερμαντικά σώματα αερίου με εξαερισμό,φούρνους πετρελαίου,καπνό τσιγάρου,θερμαντικά σώματα κεροζίνης,σόμπες ή φούρνους με ξύλο ή άνθρακα [60].
Επιπλέον, ο εξωτερικός αέρας θεωρείται επίσης κύρια πηγή εσωτερικού SO₂ [61]. Οι συγκεντρώσεις SO₂ σε εσωτερικούς χώρους είναι συχνά χαμηλότερες από τις αντίστοιχες εξωτερικές. Οι εσωτερικές εκπομπές SO₂ είναι συνήθως μικρές, λόγω της επαναδραστηριοποίησης και της εύκολης απορρόφησής του από επιφάνειες μέσα στο κτίριο. Είναι γνωστό ότι η ωριαία συγκέντρωση SO₂ σε κτίρια συχνά είναι κάτω από 20 ppb [62].

Η έκθεση του ανθρώπου στο SO₂ γίνεται αποκλειστικά μέσω εισπνοής, και μπορεί να προκαλέσει δυσλειτουργία του αναπνευστικού συστήματος.

Sulfur dioxide (SO2) is the most common gas among the group of sulfur oxides (SOx) present in the atmosphere. SO2 is primarily produced by the combustion process of fossil fuels, and combines with aerosols and PMs to form a complex group of distinct air [59]. Indoor sources of SO2 emissions include vented gas appliances, oil furnaces, tobacco smoke, kerosene heaters, and coal or wood stoves [60]. In addition, outdoor air is also regarded as a main source of indoor SO2 [61]. Indoor SO2 levels are often lower than outdoor levels. SO2 emission indoors is usually small, owing to its reactivation, which can be easily absorbed by indoor surfaces. It is known that the hourly concentration of SO2 in buildings is often below 20 ppb [62]. Human exposure to SO2, which can impair respiratory function, is only via inhalation

\cite{ijerph17238972}


\subsubsection{Aerosols}

Indoor aerosols are either primary aerosols originating from different indoor sources or secondary aerosols formed by indoor gas-to-particle reactions [79]. Moreover, outdoor particles infiltrating indoors are also likely to be a source of indoor aerosols. Secondary inorganic aerosols are PMs consisting of inorganic elements, including anthropogenic or crustal sources and water-soluble ions [6], while secondary organic aerosols (SOAs) are formed in the gas-to-particle conversion process of VOCs [80]. Additionally, carbonaceous aerosols, which comprise SOAs and elemental carbons released in incomplete combustion, are well-known species in PM2.5 [81]. Biological aerosols (bioaerosols) are a subset of atmospheric PMs comprising dispersal units (fungal spores and plant pollen), microorganisms (bacteria and archaea), or cellular materials [82]. Due to their diversity in terms of compounds and phases (gas, liquid, or solid), aerosols can be regarded as dynamic systems [83]. As such, their particle size distribution varies from the nucleation mode (<30 nm in vacuum cleaning condition) to the accumulation mode (~100 nm, indoor combustion aerosols from smoking, cooking, or incense burning), and to the fine and coarse modes (>1 µm, resuspension aerosols) [84,85]. Aerosol exposure through inhalation in the indoor environment has been linked to numerous adverse health effects, mainly in the lungs (the entrance to the human body) and other important target organs, such as the heart and brain [79].


\cite{ijerph17238972}



\subsubsection{Radon}

Οι κύριες πηγές ραδονίου στους εσωτερικούς χώρους περιλαμβάνουν τα οικοδομικά υλικά, τα αέρια του εδάφους και το νερό της βρύσης [86]. Καθώς το έδαφος περιέχει ράδιο σε ίχνη, το ράδιον είναι πιθανό να αποτελεί ένα από τα συστατικά του αερίου που γεμίζει τους πόρους του εδάφους. Όσον αφορά τις εκπομπές ραδονίου από οικοδομικά υλικά, όλα τα υλικά που περιέχουν ίχνη ραδίου μπορούν να απελευθερώσουν ράδιον. Μεταξύ των οικοδομικών υλικών, τα δομικά υλικά όπως πέτρα, σκυρόδεμα και τούβλο αποτελούν τις κύριες πηγές εκπομπής ραδονίου στους εσωτερικούς χώρους, δεδομένου ότι χρησιμοποιούνται τόνοι αυτών των υλικών στην κατασκευή κτιρίων.

Το ράδιον στους εσωτερικούς χώρους μπορεί επίσης να απελευθερωθεί μέσω της χρήσης νερού από υπόγειες πηγές που περιέχουν γρανίτη ή άλλα πετρώματα με ράδιο, και τέτοιες πηγές νερού συχνά περιέχουν συγκεντρώσεις ραδονίου άνω των 10.000 pCi/L [87]. Τέλος, ο εξωτερικός αέρας θεωρείται επίσης πηγή εισόδου ραδονίου στους εσωτερικούς χώρους [88]. Η ανθρώπινη έκθεση στο ράδιον μέσα σε κτίρια πραγματοποιείται κυρίως μέσω της διάχυσης αερίων από το έδαφος που βρίσκεται κάτω από τα κτίρια [89]. Επιδημιολογικές μελέτες έχουν δείξει ότι το ράδιον στους εσωτερικούς χώρους μπορεί να αυξήσει τον κίνδυνο καρκίνου του πνεύμονα κατά 3% έως 14%, ανάλογα με το μέσο επίπεδο ραδονίου [90].

\cite{ijerph17238972}


The primary sources of indoor radon include building materials, soil gas, and tap water [86]. As soil contains radium at trace concentrations, radon is likely to be one of the constituents in the gas filling soil pores. As for radon emissions from building materials, all materials holding trace amounts of radium can release radon. Among building materials, masonry materials (i.e., stone, concrete, and brick) are the main sources for indoor radon emission, in that tons of such materials are used in building construction. Indoor radon can be released through the usage of water from underground water sources containing granite or other radium-bearing rock, and such water sources commonly contain radon concentrations above 10,000 pCi/L [87]. Finally, outdoor air is also regarded as a source of indoor radon [88]. Human exposure to radon in buildings is incurred mainly through the permeation pathways of underlying soil gas [89]. Epidemiological studies have demonstrated that indoor radon can cause lung cancer risk increases of 3% to 14%, depending on the average radon level [90].

\cite{ijerph17238972}

\subsubsection{Pesticides}

These days, inorganic and organic pesticides have commonly been utilized as protectants for wooden building materials by impregnation or surface coating [91]. Pesticides are also used to control and prevent pests, including bacteria, fungi, insects, rodents, and other organisms [92,93]. In the indoor environment, pesticides are usually semi-volatile compounds that may exist in either gas or particulate form according to properties, such as the vapor pressure, product viscosity, and water solubility [94]. In addition, it has been indicated that carpet and textiles are likely to play the role of long-term reservoirs for organochlorine pesticides [95,96]. It is supposed that when used in carpets, textiles, and cushioned furniture, pesticides in fibers will migrate into polyurethane foam pads [97,98], and thus carpet, textiles, and cushioned furniture can reflect an integrated pesticide exposure during their lifetime. Moreover, pesticides are able to enter buildings from outdoors. Once inside, they can persist for months or years due to their protection against sunlight, extreme temperatures, rain, and other factors [97]. Dermal uptake, ingestion, and inhalation of particles or volatile compounds containing pesticides are believed to be potential exposure routes in the indoor environment [92]. Pesticide exposure is associated with adverse health risks, including (i) short-term skin and eye irritation, dizziness, headaches, and nausea; and (ii) long-term chronic impacts, such as cancer, asthma, and diabetes [99].

\cite{ijerph17238972}



\subsubsection{ Biological Pollutants}


Οι βιολογικοί ρύποι στους εσωτερικούς χώρους περιλαμβάνουν βιολογικούς αλλεργιογόνους παράγοντες (π.χ. τρίχες και σάλιο ζώων, σκόνη σπιτιού, κατσαρίδες, ακάρεα και γύρη) και μικροοργανισμούς (ιούς, μύκητες και βακτήρια) [100].

Οι βιολογικοί αλλεργιογόνοι παράγοντες, γνωστοί και ως αντιγόνα, προέρχονται από διάφορα έντομα, ζώα, ακάρεα, φυτά ή μύκητες, και προκαλούν αλλεργική αντίδραση όταν αλληλεπιδρούν με τα ειδικά αντισώματα ανοσοσφαιρίνης Ε (IgE) [101]. Οι κύριες πηγές αλλεργιογόνων στους εσωτερικούς χώρους περιλαμβάνουν κατοικίδια με τρίχωμα (π.χ. σκύλους και γάτες), ακάρεα σπιτιού, μούχλα, φυτά, κατσαρίδες και τρωκτικά [102], ενώ υπάρχουν και εξωτερικές πηγές [101].

Οι ιοί και τα βακτήρια προέρχονται συχνά από ανθρώπους και ζώα ή μεταφέρονται από αυτά. Έχει αποδειχθεί ότι η έκθεση σε βιολογικούς αλλεργιογόνους παράγοντες μπορεί να προκαλέσει ευαισθητοποίηση, αναπνευστικές λοιμώξεις, αλλεργικές αναπνευστικές παθήσεις και συριγμό [103], ενώ η έκθεση σε βακτήρια και ιούς στους εσωτερικούς χώρους πιθανόν να προκαλέσει τόσο μολυσματικές όσο και μη μολυσματικές αρνητικές επιπτώσεις στην υγεία [104].

\cite{ijerph17238972}

Biological pollutants in indoor environments include biological allergens (e.g., animal dander and cat saliva, house dust, cockroaches, mites, and pollen) and microorganisms (viruses, fungi, and bacteria) [100]. Biological allergens, known as antigens, originate from a number of insects, animals, mites, plants, or fungi, and will induce an allergic state in reacting with specific immunoglobulin E (IgE) antibodies [101]. Indoor sources of allergens mainly include furred pets (dog and cat dander), house dust mites, molds, plants, cockroaches, and rodents [102], and there are outdoor sources as well [101]. Viruses and bacteria often originate from or are carried by people and animals. It has been demonstrated that exposure to biological allergens can result in sensitization, respiratory infections, respiratory allergic diseases, and wheezing [103], while exposure to bacteria and viruses indoors is likely to cause noninfectious and infectious adverse health outcomes [104].


\cite{ijerph17238972}
